
\documentclass[11pt,a4paper,twoside]{article}
\usepackage[inner=28mm,outer=22mm,top=22mm,bottom=23mm,headheight=15pt]{geometry}
\usepackage[utf8]{inputenc}
\usepackage[T1]{fontenc}
\usepackage{lmodern}
\usepackage[provide=*,spanish]{babel}
\usepackage{microtype,graphicx,longtable,booktabs,tabularx,array,xcolor,enumitem,hyperref,fancyhdr,caption,float,pdflscape,seqsplit,cite,lastpage,pdfpages}
\usepackage[most]{tcolorbox}
\graphicspath{{../03_Modelado/Diagramas_UML/}{modelado_final/}}
\definecolor{navy}{HTML}{123A5A}
\definecolor{teal}{HTML}{2A7F85}
\definecolor{coral}{HTML}{C85A3E}
\definecolor{lightblue}{HTML}{EAF3F6}
\hypersetup{colorlinks=true,linkcolor=navy,urlcolor=teal,citecolor=navy,pdftitle={FabroGym ERS/SRS 2B v2.0},pdfauthor={Equipo FabroGym ISR-401}}
\pagestyle{fancy}
\fancyhf{}
\fancyhead[LE]{\small FabroGym -- ERS/SRS 2B}
\fancyhead[RO]{\small Entrega 4 (2B) -- v2.0}
\fancyfoot[LE,RO]{\thepage\ de \pageref{LastPage}}
\setlength{\parindent}{0pt}
\setlength{\parskip}{4pt}
\setlist{nosep,leftmargin=*}
\newtcolorbox{importantbox}{enhanced,breakable,colback=lightblue,colframe=navy,boxrule=.6pt,arc=1mm,left=6pt,right=6pt,top=5pt,bottom=5pt}
\newtcolorbox{riskbox}{enhanced,breakable,colback=white,colframe=coral,boxrule=.8pt,arc=1mm,left=6pt,right=6pt,top=5pt,bottom=5pt}
\newcommand{\estadoReq}[1]{\textsf{#1}}
\newcommand{\zenododoi}{\href{https://doi.org/10.5281/zenodo.22237884}{10.5281/zenodo.22237884}}
\begin{document}

\begin{titlepage}
\centering
{\Large\bfseries UNIVERSIDAD TÉCNICA ESTATAL DE QUEVEDO\par}
\vspace{0.25cm}
{\large Facultad de Ciencias de la Computación\par}
{\large Carrera de Ingeniería de Software\par}
\vspace{0.45cm}
{\large Asignatura: Ingeniería de Requerimientos (ISR-401)\par}
\vspace{0.55cm}
{\LARGE\bfseries Especificación de Requisitos de Software (ERS/SRS)\par}
\vspace{0.25cm}
{\Large\bfseries Entrega 4 (2B / Defensa Final) -- Versión 2.0\par}
\vspace{0.18cm}
{\large FabroGym -- sistema de gestión administrativa y operativa para un gimnasio\par}
\vspace{0.5cm}
\begin{tabularx}{\textwidth}{|>{\bfseries}p{4.0cm}|X|}\hline
Proyecto & FabroGym \\\hline
Versión & 2.0 -- entrega académica vigente con UML definitivo \\\hline
Fecha de corte documental & 17 de septiembre de 2026 \\\hline
Repositorio & \url{https://github.com/gleiston-guerrero/FabroGym_ISR401} \\\hline
Registro OSF & \url{https://osf.io/62ysc/} -- DOI \href{https://doi.org/10.17605/OSF.IO/62YSC}{10.17605/OSF.IO/62YSC} \\\hline
DOI del paquete Zenodo & \zenododoi \\\hline
Estado & Documento 2B consolidado con el conjunto UML definitivo del equipo integrado y verificado contra los 19 RF Must, la arquitectura terminal y los límites de alcance. \\\hline
\end{tabularx}
\vspace{0.35cm}
{\normalsize\bfseries Equipo y autoría del proyecto en el cierre del examen suspenso\par}
\vspace{0.1cm}
\small
\begin{tabularx}{\textwidth}{|X|X|p{4.15cm}|}\hline
\textbf{Integrante} & \textbf{Situación documental} & \textbf{ORCID}\\\hline
Mera Arias Erick Jhair & Autoría preservada; trabajo posterior a la guía verificado & 0009-0001-0068-1796\\\hline
Ponce Rivera Mery Helenmey & Autoría preservada; trabajo posterior a la guía verificado & 0009-0006-6041-9198\\\hline
Mora Duarte Alex José & Autoría preservada; aportes históricos verificables & 0009-0000-2494-2842\\\hline
Alvia Villegas Erick Adalberto & Autoría preservada; aportes históricos verificables & 0009-0001-3777-470X\\\hline
Vaca Romero David Octavio & Autoría preservada; aportes históricos verificables & 0009-0000-4457-3095\\\hline
\end{tabularx}
\vspace{0.12cm}
\footnotesize \textit{Los cinco integrantes conservan la autoría de sus aportes verificables. Para el corte posterior a la guía, el historial distingue la actividad efectivamente versionada en ese periodo de los aportes históricos previos. La interpretación vigente se documenta en 04\_Trazabilidad/ACTA\_RECONOCIMIENTO\_AUTORIA\_EQUIPO.pdf y 10\_Autoria/EQUIPO\_EXAMEN\_FINAL.md; la solicitud del 15/09/2026 se conserva únicamente como antecedente documental.}
\vfill
\normalsize Quevedo, Ecuador -- 2026
\end{titlepage}
\tableofcontents
\listoffigures
\listoftables
\newpage

\section{Control de versión y estado}
Para la Entrega 4 (2B / Defensa Final), la única versión académica vigente es la \textbf{2.0}. Las denominaciones 2.1 y 2.2 se conservan exclusivamente como revisiones internas históricas empleadas durante la consolidación; no constituyen entregables vigentes distintos. La versión 2.0 integra 54 diagramas UML, 25 RF, 23 RNF y 4 restricciones de diseño. El componente inteligente continúa propuesto y no se atribuyen pruebas o capacidades no ejecutadas. El paquete de datos fue publicado en Zenodo como versión 2.0.0 con DOI específico \zenododoi.

\begin{tabularx}{\textwidth}{p{2cm}p{2.4cm}X}\toprule
Versión & Fecha & Cambio principal \\ \midrule
2.2 (interna, no vigente) & 2026-09-01 & Revisión interna de integración del conjunto UML definitivo; queda subordinada a la versión académica 2.0.\\
2.1 (interna, no vigente) & 2026-09-01 & Ampliación acumulativa: conserva el cierre 2B y reincorpora 64 páginas de detalle técnico heredado (mockups, historias/casos de uso, UML y catálogo de interfaces), claramente marcado como anexo histórico y subordinado a la versión académica vigente 2.0.\\
2.0 & 2026-09-15 & Cierre 2B vigente para el examen suspenso: 25 RF, 23 RNF, 4 restricciones de diseño, especificación textual completa de los 19 casos de uso Must y 162 filas de trazabilidad (97 trazas preexistentes migradas, 8 planes de verificación del componente inteligente y 57 trazas de flujo CU); integración de 16 sesiones, encuesta de 70 respuestas, análisis reproducible y OSF v1.4.\\
1.5.8 & 2026-08-02 & Consolidación 2A con diez entrevistas, codificación temática inicial, 44 identificadores y 70 mockups.\\
1.5.4 & 2026-08-02 & Compatibilidad de compilación y consolidación visual.\\
0.2.0 & 2026-07-12 & Entrega anterior utilizada como referencia estructural.\\
\bottomrule\end{tabularx}

\begin{importantbox}
\textbf{Regla de integridad.} Esta ERS diferencia lo \textbf{especificado}, lo \textbf{prototipado} y lo \textbf{propuesto}. El componente de recomendación basado en IA/AA continúa \textbf{propuesto}; los RNF de explicabilidad no convierten ese componente en una funcionalidad implementada.
\end{importantbox}

\section{Introducción}
\subsection{Propósito}
Definir de forma necesaria, singular, verificable y trazable los requisitos de FabroGym para el cierre de ISR-401, enlazando evidencia de campo, decisiones de privacidad, casos de uso, historias, criterios de aceptación, arquitectura del MVP y resultados del estudio de explicabilidad. La redacción sigue los principios de ISO/IEC/IEEE 29148:2018 y la operacionalización de calidad se apoya en ISO/IEC 25010:2023 \cite{iso29148,iso25010,pohl2010,swebok2024}.

\subsection{Alcance del producto}
FabroGym cubre autenticación no biométrica, control por roles, clientes con datos mínimos, membresías, pagos, asistencia, productos, inventario, ventas, cierre de caja, novedades, horarios de instructor, rutinas y reportes. Se excluyen facturación electrónica integrada, pasarelas bancarias, nómina, marketing automatizado, reconocimiento biométrico, salud, diagnóstico, nutrición, peso o medidas corporales reales e IA/AA operativa.

\subsection{Alcance del componente de recomendación}
La idea de recomendación de rutinas se mantiene como evolución propuesta. El trabajo empírico 2B se centra en requisitos de explicabilidad que deberán cumplirse si ese componente se implementa en una versión futura. Por tanto, sus RNF se especifican y validan conceptualmente, pero no se presentan como comportamiento disponible en el MVP.

\subsection{Glosario}
\begin{tabularx}{\textwidth}{p{2.7cm}X}\toprule
Término & Definición \\ \midrule
RF & Requisito funcional.\\
RNF & Requisito no funcional / requisito de calidad.\\
RD & Restricción de diseño o alcance.\\
RNF de explicabilidad & Requisito no funcional derivado del Enfoque 3; conserva Explicabilidad como característica.\\
EV & Identificador de evidencia.\\
CU / HU / CA & Caso de uso / historia de usuario / criterio de aceptación.\\
WALK-TEC / WALK-NTEC & Walkthrough técnico / no técnico; se preserva la técnica original.\\
MC-P01..03 & Seudónimos de participantes de member checking documental.\\
MVP & Producto mínimo viable demostrativo con datos sintéticos.\\
\bottomrule\end{tabularx}

\section{Fuentes, método y evidencia de requisitos}
\subsection{Fuentes de verdad}
La prioridad de decisión es: guía/rúbrica 2B y paquete ético; evidencia primaria real; registro/protocolo OSF; repositorio; ERS heredada; estándares y literatura. Ante una contradicción, la evidencia primaria prevalece sobre inferencias y un artefacto terminal prevalece sobre una versión histórica.

\subsection{Evidencia acumulada y nomenclatura}
La evidencia audiovisual se compone de diez entrevistas iniciales (ENTR-01 a ENTR-10) y seis walkthroughs (WALK-NTEC-01 a WALK-NTEC-03; WALK-TEC-01 a WALK-TEC-03). Los seis walkthroughs forman parte del acumulado de 16 sesiones aceptado para el cierre, sin renombrarlos como entrevistas. La ficha técnica v3.1 reporta 16 videos con duración total de \textbf{06:18:08} y 16 audios con duración total de \textbf{06:18:14}; para el mínimo temporal se cuenta una duración por sesión y no se duplican audio y video.

\begin{table}[H]\centering\small
\begin{tabularx}{\textwidth}{>{\raggedright\arraybackslash}p{3.1cm} >{\centering\arraybackslash}p{1.8cm} >{\raggedright\arraybackslash}p{5.1cm} X}\toprule
Fuente & N & Desagregación & Uso en requisitos \\ \midrule
Entrevistas ENTR & 10 & propietarios/administración, recepción e instructores & necesidades operativas, privacidad, reglas de negocio, RF y RNF base.\\
Walkthroughs WALK & 6 & 3 técnicos + 3 no técnicos & validación formativa, codificación terminal y explicabilidad.\\
Encuesta clientes & 70 & sin variable técnico/no técnico & evidencia descriptiva del dominio; no se convierte en Likert de explicabilidad.\\
Member checking & 3 participantes & MC-P01, MC-P02, MC-P03 & 12 decisiones sobre 4 candidatos RNF de explicabilidad.\\
\bottomrule\end{tabularx}
\caption{Corpus empírico utilizado en el cierre 2B.}\label{tab:corpus}
\end{table}

\subsection{Codificación y explicabilidad}
El corpus de walkthroughs contiene 76 fragmentos codificados: 49 técnicos y 27 no técnicos, con 37 códigos normalizados y 18 categorías. Nueve fragmentos se clasificaron como pertinentes para explicabilidad (7 técnicos y 2 no técnicos). De ellos se derivaron cuatro candidatos RNF, revisados documentalmente el 29 de agosto de 2026 por MC-P01, MC-P02 y MC-P03. El registro contiene 12 decisiones: 4 confirmaciones, 8 ajustes y 0 rechazos. El proceso de member checking sí ocurrió; no existe grabación audiovisual de esa actividad y no se afirma lo contrario \cite{birt2016}.

\subsection{Saturación, tamaño del efecto y límite estadístico}
A nivel de códigos normalizados, las últimas tres sesiones incorporan un promedio de 2,333 códigos nuevos sobre 37 acumulados, equivalente a 6,306\%, por lo que no se declara cumplimiento estricto del umbral de 5\%. A nivel de categorías axiales el valor complementario es 1,852\%, lo que muestra estabilización conceptual. Esta diferencia se conserva como limitación transparente y no se reetiqueta para forzar un resultado \cite{hennink2017}.

Para el contraste técnico/no técnico se corrigió la unidad de análisis: las 18 categorías temáticas no se tratan como 18 observaciones independientes. La unidad es cada sesión WALK, con \texttt{n\_unidades=6} (3 técnicas y 3 no técnicas). Sobre la proporción de fragmentos pertinentes a explicabilidad por sesión, el delta de Cliff es 0,555556 y su IC95\% bootstrap exacto es [-0,333333; 1,000000] \cite{cliff1993}. La salida terminal marca \texttt{interpretable=NO} para inferencia poblacional debido al reducido número de sesiones independientes y a la amplitud del intervalo; el cálculo se conserva únicamente como estimación descriptivo-exploratoria del caso. El análisis secundario sobre conteos brutos produce delta=0,777778 e IC95\% [0,333333; 1,000000], pero tampoco se interpreta inferencialmente porque el número total de fragmentos difiere entre sesiones.

\begin{table}[H]\centering\small
\begin{tabularx}{\textwidth}{p{4.8cm}c c p{3.4cm}c}\toprule
Métrica & $n$ unidades & Delta de Cliff & IC95\% exacto & Interpretable \\ \midrule
Proporción pertinente por sesión & 6 (3+3) & 0,555556 & [-0,333333; 1,000000] & NO \\
Conteo pertinente por sesión (sensibilidad) & 6 (3+3) & 0,777778 & [0,333333; 1,000000] & NO \\ \bottomrule
\end{tabularx}
\caption{Auditoría del tamaño del efecto técnico/no técnico por unidades independientes. \texttt{Interpretable=NO} se refiere exclusivamente a inferencia poblacional.}\label{tab:efecto-perfiles}
\end{table}

La encuesta de 70 respuestas no contiene perfil técnico/no técnico ni ítems Likert de explicabilidad por dimensión. En consecuencia, no se usa para el contraste entre perfiles ni se calculan Mann--Whitney, Fleiss $\kappa$ o una satisfacción Likert de explicabilidad con datos inexistentes. Los índices ordinales del cuestionario permanecen descriptivos y separados del análisis de las seis sesiones WALK \cite{molleri2020,cohen1988}.

\section{Descripción general del sistema}
\subsection{Perspectiva y límites}
FabroGym es un sistema web demostrativo y modular. El MVP utiliza HTML5, CSS3 y JavaScript y trabaja con datos sintéticos en almacenamiento local. La versión de demostración puede ejecutarse abriendo el sitio directamente, con un servidor HTTP local o mediante un contenedor web. El sistema no integra proveedores bancarios, facturación electrónica ni mensajería externa.

\begin{figure}[H]\centering\includegraphics[width=.94\textwidth]{figuras_2B/fig_contexto_2B.pdf}\caption{Contexto del sistema y límites de alcance 2B.}\label{fig:contexto}\end{figure}

\subsection{Usuarios y responsabilidades}
\begin{tabularx}{\textwidth}{p{3.4cm}X}\toprule
Rol & Responsabilidades \\ \midrule
Administrador & Configuración, catálogo de planes, productos, permisos, auditoría, cierres y reportes.\\
Recepción & Clientes, membresías, pagos, asistencia, ventas, novedades y operación cotidiana.\\
Instructor & Rutinas, versiones, seguimiento operativo y disponibilidad, sin datos de salud reales.\\
Cliente & Beneficiario de la operación; su autoservicio no es esencial al MVP actual.\\
\bottomrule\end{tabularx}

\subsection{Supuestos y dependencias}
Se asumen navegadores modernos, credenciales demostrativas y datos sintéticos. La demostración no equivale a un despliegue productivo. La producción requeriría un backend persistente, cifrado en tránsito, controles de acceso verificables, respaldo y procedimiento formal de derechos del titular.

\section{Catálogo terminal de requisitos}
El catálogo canónico elimina IDs duplicados de versiones previas. La migración determinista conserva 25 RF, 19 RNF preexistentes y 4 RD sin pérdida ni duplicación. Además se incorporan RNF-20 a RNF-23 como propuestas normativas para recomendación, equidad, monitoreo y riesgo. La ERS contiene \textbf{52 requisitos terminales}: 25 RF, 23 RNF y 4 RD. El mapa completo está en \path{04_Trazabilidad/mapeo_ids_2B.csv}.

\subsection{Requisitos funcionales}
\subsubsection{RF-01 -- Garantizar tratamiento lícito y minimización de datos}
\begin{tabularx}{\textwidth}{|>{\bfseries}p{3.2cm}|X|}\hline
Enunciado & El sistema deberá permitir iniciar y cerrar sesión mediante credenciales individuales no biométricas. \\ \hline
Actor/rol & Personal autorizado \\ \hline
Prioridad / Kano & Must / Básico \\ \hline
Evidencia & EV-ENTR02-07 \\ \hline
CU / HU / CA & CU-01 / HU-01 / CA-01 \\ \hline
Componente & Seguridad \\ \hline
Mockups & MU-01 \\ \hline
Estado & ESPECIFICADO; cobertura MVP se verifica en C3 \\ \hline
\end{tabularx}
\textbf{Criterio de aceptación terminal.} Dadas credenciales de demostración válidas, cuando el usuario inicia sesión, entonces se crea una sesión del rol correspondiente; con credenciales inválidas se deniega el acceso sin revelar qué campo falló.
\subsubsection{RF-02 -- Controlar acceso por roles y permisos}
\begin{tabularx}{\textwidth}{|>{\bfseries}p{3.2cm}|X|}\hline
Enunciado & El sistema deberá autorizar cada operación según el rol Administrador, Recepción o Instructor y denegar por defecto las no asignadas. \\ \hline
Actor/rol & Administrador \\ \hline
Prioridad / Kano & Must / Básico \\ \hline
Evidencia & EV-ENTR02-07;EV-ENTR03-06;EV-ENTR08-04 \\ \hline
CU / HU / CA & CU-02 / HU-02 / CA-02 \\ \hline
Componente & Seguridad \\ \hline
Mockups & MU-22;MU-32 \\ \hline
Estado & ESPECIFICADO; cobertura MVP se verifica en C3 \\ \hline
\end{tabularx}
\textbf{Criterio de aceptación terminal.} Dado un usuario autenticado, cuando intenta una operación fuera de su rol, entonces el sistema la deniega; cuando la operación está autorizada, permite continuar.
\subsubsection{RF-03 -- Asegurar trazabilidad de eventos críticos}
\begin{tabularx}{\textwidth}{|>{\bfseries}p{3.2cm}|X|}\hline
Enunciado & El sistema deberá registrar autor, fecha, acción y valores anterior/nuevo para pagos, precios, inventario, membresías y permisos. \\ \hline
Actor/rol & Administrador \\ \hline
Prioridad / Kano & Should / Desempeño \\ \hline
Evidencia & EV-ENTR02-04;EV-ENTR02-07;EV-ENTR03-06 \\ \hline
CU / HU / CA &  /  / CA-N/A \\ \hline
Componente & Auditoría \\ \hline
Mockups & MU-25 \\ \hline
Estado & ESPECIFICADO; cobertura MVP se verifica en C3 \\ \hline
\end{tabularx}
\subsubsection{RF-04 -- Registrar datos mínimos de cliente}
\begin{tabularx}{\textwidth}{|>{\bfseries}p{3.2cm}|X|}\hline
Enunciado & El sistema deberá registrar un cliente con código interno, nombre de uso y un medio de contacto, sin salud, biometría, dirección ni imagen corporal. \\ \hline
Actor/rol & Recepción \\ \hline
Prioridad / Kano & Must / Desempeño \\ \hline
Evidencia & EV-ENTR01-02;EV-ENTR02-02;EV-ENTR03-02;EV-ENTR08-01 \\ \hline
CU / HU / CA & CU-03 / HU-03 / CA-03 \\ \hline
Componente & Clientes \\ \hline
Mockups & MU-12;MU-37 \\ \hline
Estado & ESPECIFICADO; cobertura MVP se verifica en C3 \\ \hline
\end{tabularx}
\textbf{Criterio de aceptación terminal.} Dado un operador autorizado, cuando registra código interno, nombre de uso y contacto permitido, entonces se crea un cliente único sin solicitar salud, biometría, dirección ni imagen corporal.
\subsubsection{RF-05 -- Consultar ficha de cliente}
\begin{tabularx}{\textwidth}{|>{\bfseries}p{3.2cm}|X|}\hline
Enunciado & El sistema deberá buscar clientes por código, nombre normalizado o contacto y mostrar membresía, último pago, asistencia y novedades según permisos. \\ \hline
Actor/rol & Recepción \\ \hline
Prioridad / Kano & Must / Desempeño \\ \hline
Evidencia & EV-ENTR02-03;EV-ENTR02-08;EV-ENTR08-05 \\ \hline
CU / HU / CA & CU-04 / HU-04 / CA-04 \\ \hline
Componente & Clientes \\ \hline
Mockups & MU-11;MU-36 \\ \hline
Estado & ESPECIFICADO; cobertura MVP se verifica en C3 \\ \hline
\end{tabularx}
\textbf{Criterio de aceptación terminal.} Dado un criterio de búsqueda válido, cuando se consulta un cliente, entonces se muestra su ficha operativa permitida sin exponer información fuera del rol.
\subsubsection{RF-06 -- Actualizar y reactivar cliente}
\begin{tabularx}{\textwidth}{|>{\bfseries}p{3.2cm}|X|}\hline
Enunciado & El sistema deberá actualizar o reactivar un registro existente y advertir coincidencias antes de crear otro cliente. \\ \hline
Actor/rol & Recepción \\ \hline
Prioridad / Kano & Must / Desempeño \\ \hline
Evidencia & EV-ENTR02-02;EV-ENTR03-02;EV-ENTR04-01;EV-ENTR05-02 \\ \hline
CU / HU / CA & CU-05 / HU-05 / CA-05 \\ \hline
Componente & Clientes \\ \hline
Mockups & MU-11;MU-36 \\ \hline
Estado & ESPECIFICADO; cobertura MVP se verifica en C3 \\ \hline
\end{tabularx}
\textbf{Criterio de aceptación terminal.} Dado un cliente existente, cuando se actualiza o reactiva, entonces se conserva el mismo identificador y se advierten coincidencias antes de crear otro registro.
\subsubsection{RF-07 -- Configurar catálogo de planes}
\begin{tabularx}{\textwidth}{|>{\bfseries}p{3.2cm}|X|}\hline
Enunciado & El sistema deberá administrar nombre, duración, precio, vigencia y estado de planes y promociones. \\ \hline
Actor/rol & Administrador \\ \hline
Prioridad / Kano & Must / Desempeño \\ \hline
Evidencia & EV-ENTR02-01;EV-ENTR04-01 \\ \hline
CU / HU / CA & CU-06 / HU-06 / CA-06 \\ \hline
Componente & Membresías \\ \hline
Mockups & MU-68 \\ \hline
Estado & ESPECIFICADO; cobertura MVP se verifica en C3 \\ \hline
\end{tabularx}
\textbf{Criterio de aceptación terminal.} Dado un administrador, cuando crea o modifica un plan, entonces quedan registrados nombre, duración, precio, vigencia y estado y el catálogo refleja el cambio.
\subsubsection{RF-08 -- Activar y renovar membresía}
\begin{tabularx}{\textwidth}{|>{\bfseries}p{3.2cm}|X|}\hline
Enunciado & El sistema deberá activar o renovar una membresía calculando inicio y vencimiento desde un plan vigente y un pago confirmado. \\ \hline
Actor/rol & Recepción \\ \hline
Prioridad / Kano & Must / Básico \\ \hline
Evidencia & EV-ENTR01-02;EV-ENTR02-04;EV-ENTR03-01;EV-ENTR04-03 \\ \hline
CU / HU / CA & CU-07 / HU-07 / CA-07 \\ \hline
Componente & Membresías \\ \hline
Mockups & MU-14;MU-39 \\ \hline
Estado & ESPECIFICADO; cobertura MVP se verifica en C3 \\ \hline
\end{tabularx}
\textbf{Criterio de aceptación terminal.} Dado un plan vigente y un pago confirmado, cuando se activa o renueva una membresía, entonces el inicio y vencimiento se calculan y persisten de forma consistente.
\subsubsection{RF-09 -- Consultar vigencia}
\begin{tabularx}{\textwidth}{|>{\bfseries}p{3.2cm}|X|}\hline
Enunciado & El sistema deberá mostrar estado, plan, fecha de inicio y vencimiento de la membresía. \\ \hline
Actor/rol & Personal autorizado \\ \hline
Prioridad / Kano & Must / Básico \\ \hline
Evidencia & EV-ENTR01-03;EV-ENTR02-03;EV-ENTR03-03;EV-ENTR08-02 \\ \hline
CU / HU / CA & CU-08 / HU-08 / CA-08 \\ \hline
Componente & Membresías \\ \hline
Mockups & MU-13;MU-38;MU-63(F) \\ \hline
Estado & ESPECIFICADO; cobertura MVP se verifica en C3 \\ \hline
\end{tabularx}
\textbf{Criterio de aceptación terminal.} Dado un cliente localizado, cuando se consulta su membresía, entonces se muestran estado, plan, inicio y vencimiento.
\subsubsection{RF-10 -- Alertar vencimientos}
\begin{tabularx}{\textwidth}{|>{\bfseries}p{3.2cm}|X|}\hline
Enunciado & El sistema deberá listar diariamente membresías que vencen en los próximos tres días y las vencidas, sin enviar mensajes externos automáticamente. \\ \hline
Actor/rol & Recepción \\ \hline
Prioridad / Kano & Must / Desempeño \\ \hline
Evidencia & EV-ENTR01-03;EV-ENTR05-02;EV-ENTR07-02 \\ \hline
CU / HU / CA & CU-09 / HU-09 / CA-09 \\ \hline
Componente & Membresías \\ \hline
Mockups & MU-10;MU-13;MU-38 \\ \hline
Estado & ESPECIFICADO; cobertura MVP se verifica en C3 \\ \hline
\end{tabularx}
\textbf{Criterio de aceptación terminal.} Dada la fecha de corte, cuando se consulta la alerta, entonces se listan las membresías vencidas y las que vencen dentro de tres días sin enviar mensajes externos automáticamente.
\subsubsection{RF-11 -- Registrar pago y comprobante}
\begin{tabularx}{\textwidth}{|>{\bfseries}p{3.2cm}|X|}\hline
Enunciado & El sistema deberá registrar monto, concepto, medio, referencia y fecha de un pago, y generar un comprobante interno. \\ \hline
Actor/rol & Recepción \\ \hline
Prioridad / Kano & Must / Básico \\ \hline
Evidencia & EV-ENTR01-02;EV-ENTR02-04;EV-ENTR03-04;EV-ENTR07-02;EV-ENTR08-03 \\ \hline
CU / HU / CA & CU-10 / HU-10 / CA-10 \\ \hline
Componente & Pagos \\ \hline
Mockups & MU-15;MU-16;MU-40;MU-41 \\ \hline
Estado & ESPECIFICADO; cobertura MVP se verifica en C3 \\ \hline
\end{tabularx}
\textbf{Criterio de aceptación terminal.} Dado un pago válido, cuando se registra, entonces se conservan monto, concepto, medio, referencia y fecha y se genera un comprobante interno.
\subsubsection{RF-12 -- Resolver pagos pendientes}
\begin{tabularx}{\textwidth}{|>{\bfseries}p{3.2cm}|X|}\hline
Enunciado & El sistema deberá permitir confirmar, rechazar o corregir un pago pendiente conservando motivo y auditoría. \\ \hline
Actor/rol & Administrador \\ \hline
Prioridad / Kano & Should / Desempeño \\ \hline
Evidencia & EV-ENTR02-04;EV-ENTR03-04;EV-ENTR07-07 \\ \hline
CU / HU / CA &  /  / CA-N/A \\ \hline
Componente & Pagos \\ \hline
Mockups & MU-15;MU-40 \\ \hline
Estado & ESPECIFICADO; cobertura MVP se verifica en C3 \\ \hline
\end{tabularx}
\subsubsection{RF-13 -- Registrar asistencia}
\begin{tabularx}{\textwidth}{|>{\bfseries}p{3.2cm}|X|}\hline
Enunciado & El sistema deberá registrar una asistencia no biométrica en una acción tras validar identidad operativa y vigencia, o documentar una excepción autorizada. \\ \hline
Actor/rol & Recepción \\ \hline
Prioridad / Kano & Must / Básico \\ \hline
Evidencia & EV-ENTR02-03;EV-ENTR03-03;EV-ENTR04-02;EV-ENTR07-07;EV-ENTR08-02 \\ \hline
CU / HU / CA & CU-11 / HU-11 / CA-11 \\ \hline
Componente & Asistencia \\ \hline
Mockups & MU-20;MU-30;MU-46;MU-47 \\ \hline
Estado & ESPECIFICADO; cobertura MVP se verifica en C3 \\ \hline
\end{tabularx}
\textbf{Criterio de aceptación terminal.} Dado un cliente con vigencia válida, cuando se registra asistencia, entonces se crea un único evento no biométrico; si requiere excepción, se documenta de forma autorizada.
\subsubsection{RF-14 -- Consultar asistencia}
\begin{tabularx}{\textwidth}{|>{\bfseries}p{3.2cm}|X|}\hline
Enunciado & El sistema deberá filtrar asistencias por cliente, fecha y turno, y mostrar conteos operativos. \\ \hline
Actor/rol & Personal autorizado \\ \hline
Prioridad / Kano & Must / Desempeño \\ \hline
Evidencia & EV-ENTR03-01;EV-ENTR07-01;EV-ENTR08-05;EV-ENTR09-06 \\ \hline
CU / HU / CA & CU-12 / HU-12 / CA-12 \\ \hline
Componente & Asistencia \\ \hline
Mockups & MU-20;MU-46;MU-56;MU-65(F) \\ \hline
Estado & ESPECIFICADO; cobertura MVP se verifica en C3 \\ \hline
\end{tabularx}
\textbf{Criterio de aceptación terminal.} Dado un usuario autorizado, cuando filtra asistencias por cliente, fecha o turno, entonces se muestran el detalle y los conteos correspondientes.
\subsubsection{RF-15 -- Administrar productos}
\begin{tabularx}{\textwidth}{|>{\bfseries}p{3.2cm}|X|}\hline
Enunciado & El sistema deberá registrar productos con código, nombre, precio, unidad, stock mínimo y estado. \\ \hline
Actor/rol & Administrador \\ \hline
Prioridad / Kano & Must / Desempeño \\ \hline
Evidencia & EV-ENTR01-01;EV-ENTR02-05 \\ \hline
CU / HU / CA & CU-13 / HU-13 / CA-13 \\ \hline
Componente & Inventario \\ \hline
Mockups & MU-19;MU-28;MU-44;MU-45 \\ \hline
Estado & ESPECIFICADO; cobertura MVP se verifica en C3 \\ \hline
\end{tabularx}
\textbf{Criterio de aceptación terminal.} Dado un administrador, cuando registra un producto, entonces se almacenan código, nombre, precio, unidad, stock mínimo y estado.
\subsubsection{RF-16 -- Registrar movimientos de stock}
\begin{tabularx}{\textwidth}{|>{\bfseries}p{3.2cm}|X|}\hline
Enunciado & El sistema deberá registrar entradas y ajustes con cantidad, motivo, responsable y saldo resultante. \\ \hline
Actor/rol & Personal autorizado \\ \hline
Prioridad / Kano & Must / Desempeño \\ \hline
Evidencia & EV-ENTR01-04;EV-ENTR04-04;EV-ENTR07-03 \\ \hline
CU / HU / CA & CU-14 / HU-14 / CA-14 \\ \hline
Componente & Inventario \\ \hline
Mockups & MU-69 \\ \hline
Estado & ESPECIFICADO; cobertura MVP se verifica en C3 \\ \hline
\end{tabularx}
\textbf{Criterio de aceptación terminal.} Dado un producto existente, cuando se registra una entrada o ajuste, entonces se guardan cantidad, motivo, responsable y saldo resultante.
\subsubsection{RF-17 -- Registrar venta}
\begin{tabularx}{\textwidth}{|>{\bfseries}p{3.2cm}|X|}\hline
Enunciado & El sistema deberá registrar una venta con ítems y medio de pago, y descontar existencias en la misma transacción. \\ \hline
Actor/rol & Recepción \\ \hline
Prioridad / Kano & Must / Básico \\ \hline
Evidencia & EV-ENTR02-05;EV-ENTR03-01;EV-ENTR04-04;EV-ENTR07-03;EV-ENTR08-03 \\ \hline
CU / HU / CA & CU-15 / HU-15 / CA-15 \\ \hline
Componente & Ventas \\ \hline
Mockups & MU-18;MU-27;MU-42;MU-43 \\ \hline
Estado & ESPECIFICADO; cobertura MVP se verifica en C3 \\ \hline
\end{tabularx}
\textbf{Criterio de aceptación terminal.} Dados ítems con stock suficiente, cuando se registra una venta, entonces la venta y el descuento de existencias ocurren en la misma operación demostrativa.
\subsubsection{RF-18 -- Alertar stock y rotación}
\begin{tabularx}{\textwidth}{|>{\bfseries}p{3.2cm}|X|}\hline
Enunciado & El sistema deberá listar productos en o bajo stock mínimo y resumir entradas, salidas y ventas por período. \\ \hline
Actor/rol & Administrador \\ \hline
Prioridad / Kano & Should / Desempeño \\ \hline
Evidencia & EV-ENTR01-04;EV-ENTR02-05;EV-ENTR05-04;EV-ENTR06-04 \\ \hline
CU / HU / CA &  /  / CA-N/A \\ \hline
Componente & Inventario \\ \hline
Mockups & MU-10;MU-19;MU-44 \\ \hline
Estado & ESPECIFICADO; cobertura MVP se verifica en C3 \\ \hline
\end{tabularx}
\subsubsection{RF-19 -- Cerrar turno y conciliar caja}
\begin{tabularx}{\textwidth}{|>{\bfseries}p{3.2cm}|X|}\hline
Enunciado & El sistema deberá calcular pagos y ventas por medio, comparar valores declarados y registrar diferencias y entrega de turno. \\ \hline
Actor/rol & Recepción \\ \hline
Prioridad / Kano & Must / Desempeño \\ \hline
Evidencia & EV-ENTR02-01;EV-ENTR02-06;EV-ENTR03-05;EV-ENTR04-05;EV-ENTR07-03;EV-ENTR08-04 \\ \hline
CU / HU / CA & CU-16 / HU-16 / CA-16 \\ \hline
Componente & Caja \\ \hline
Mockups & MU-70 \\ \hline
Estado & ESPECIFICADO; cobertura MVP se verifica en C3 \\ \hline
\end{tabularx}
\textbf{Criterio de aceptación terminal.} Dado un turno abierto, cuando se ejecuta el cierre, entonces se calculan pagos y ventas por medio, se compara el valor declarado y se registra la diferencia.
\subsubsection{RF-20 -- Gestionar novedades internas}
\begin{tabularx}{\textwidth}{|>{\bfseries}p{3.2cm}|X|}\hline
Enunciado & El sistema deberá crear una novedad con categoría, detalle mínimo, responsable, vencimiento y estados pendiente, en proceso, resuelta o cancelada. \\ \hline
Actor/rol & Personal autorizado \\ \hline
Prioridad / Kano & Must / Desempeño \\ \hline
Evidencia & EV-ENTR02-01;EV-ENTR02-08;EV-ENTR03-07;EV-ENTR04-06;EV-ENTR05-01 \\ \hline
CU / HU / CA & CU-17 / HU-17 / CA-17 \\ \hline
Componente & Novedades \\ \hline
Mockups & MU-21;MU-31;MU-50;MU-51;MU-61 \\ \hline
Estado & ESPECIFICADO; cobertura MVP se verifica en C3 \\ \hline
\end{tabularx}
\textbf{Criterio de aceptación terminal.} Dado un usuario autorizado, cuando crea una novedad, entonces se guardan categoría, detalle, responsable, vencimiento y estado y se conserva el historial de estado.
\subsubsection{RF-21 -- Gestionar horario y disponibilidad de instructor}
\begin{tabularx}{\textwidth}{|>{\bfseries}p{3.2cm}|X|}\hline
Enunciado & El sistema deberá registrar turnos, capacidad, disponibilidad y novedades operativas del instructor, sin datos de salud. \\ \hline
Actor/rol & Administrador \\ \hline
Prioridad / Kano & Should / Desempeño \\ \hline
Evidencia & EV-ENTR07-06;EV-ENTR09-03;EV-ENTR10-02 \\ \hline
CU / HU / CA &  /  / CA-N/A \\ \hline
Componente & Instructores \\ \hline
Mockups & MU-23;MU-33;MU-48;MU-49;MU-58 \\ \hline
Estado & ESPECIFICADO; cobertura MVP se verifica en C3 \\ \hline
\end{tabularx}
\subsubsection{RF-22 -- Administrar rutinas}
\begin{tabularx}{\textwidth}{|>{\bfseries}p{3.2cm}|X|}\hline
Enunciado & El sistema deberá crear una rutina con objetivo operativo, ejercicios, series, repeticiones, descanso y días, y asignarla a un cliente sintético o autorizado. \\ \hline
Actor/rol & Instructor \\ \hline
Prioridad / Kano & Must / Básico \\ \hline
Evidencia & EV-ENTR01-05;EV-ENTR05-04;EV-ENTR09-01;EV-ENTR09-02;EV-ENTR09-03 \\ \hline
CU / HU / CA & CU-18 / HU-18 / CA-18 \\ \hline
Componente & Rutinas \\ \hline
Mockups & MU-17;MU-29;MU-54;MU-55;MU-64(F) \\ \hline
Estado & ESPECIFICADO; cobertura MVP se verifica en C3 \\ \hline
\end{tabularx}
\textbf{Criterio de aceptación terminal.} Dado un instructor autorizado, cuando crea y asigna una rutina, entonces se registran objetivo, ejercicios, series, repeticiones, descanso y días sin datos de salud reales.
\subsubsection{RF-23 -- Registrar seguimiento}
\begin{tabularx}{\textwidth}{|>{\bfseries}p{3.2cm}|X|}\hline
Enunciado & El sistema deberá crear una nueva versión al modificar ejercicios o parámetros y conservar las versiones anteriores en solo lectura. \\ \hline
Actor/rol & Instructor \\ \hline
Prioridad / Kano & Must / Desempeño \\ \hline
Evidencia & EV-ENTR01-05;EV-ENTR09-01;EV-ENTR09-04;EV-ENTR10-01;EV-ENTR10-03 \\ \hline
CU / HU / CA & CU-19 / HU-19 / CA-19 \\ \hline
Componente & Rutinas \\ \hline
Mockups & MU-17;MU-54 \\ \hline
Estado & ESPECIFICADO; cobertura MVP se verifica en C3 \\ \hline
\end{tabularx}
\textbf{Criterio de aceptación terminal.} Dada una rutina existente, cuando se modifica un ejercicio o parámetro, entonces se crea una nueva versión y las versiones anteriores permanecen en solo lectura.
\subsubsection{RF-24 -- Registrar seguimiento operativo}
\begin{tabularx}{\textwidth}{|>{\bfseries}p{3.2cm}|X|}\hline
Enunciado & El sistema deberá registrar fecha, cumplimiento de ejercicios, repeticiones y observación operativa, excluyendo diagnósticos, lesiones, peso y medidas reales. \\ \hline
Actor/rol & Instructor \\ \hline
Prioridad / Kano & Should / Desempeño \\ \hline
Evidencia & EV-ENTR09-04;EV-ENTR09-05;EV-ENTR09-06;EV-ENTR10-03;EV-ENTR10-04 \\ \hline
CU / HU / CA &  /  / CA-N/A \\ \hline
Componente & Rutinas \\ \hline
Mockups & MU-59;MU-60 \\ \hline
Estado & ESPECIFICADO; cobertura MVP se verifica en C3 \\ \hline
\end{tabularx}
\subsubsection{RF-25 -- Consultar panel y reportes}
\begin{tabularx}{\textwidth}{|>{\bfseries}p{3.2cm}|X|}\hline
Enunciado & El sistema deberá mostrar indicadores filtrables de membresías, pagos, asistencias, ventas, caja, inventario y novedades, con exportación CSV. \\ \hline
Actor/rol & Administrador \\ \hline
Prioridad / Kano & Should / Desempeño \\ \hline
Evidencia & EV-ENTR01-05;EV-ENTR05-05;EV-ENTR06-04;EV-ENTR07-04;EV-ENTR08-05;EV-ENTR09-05 \\ \hline
CU / HU / CA &  /  / CA-N/A \\ \hline
Componente & Reportes \\ \hline
Mockups & MU-10;MU-24;MU-34 \\ \hline
Estado & ESPECIFICADO; cobertura MVP se verifica en C3 \\ \hline
\end{tabularx}

\subsection{Requisitos no funcionales de calidad}
Los RNF siguientes se redactan con condición de medición y umbral verificable. El estado ``especificado'' no implica que la prueba correspondiente ya haya sido ejecutada.
\begin{landscape}
\scriptsize
\begin{longtable}{p{2.4cm}p{3.2cm}p{10.0cm}p{4.0cm}}
\toprule ID & Característica & Requisito / métrica / umbral & Fuente \\ \midrule\endfirsthead
\toprule ID & Característica & Requisito / métrica / umbral & Fuente \\ \midrule\endhead
RNF-01 & Adecuación funcional & En cada versión candidata, el 100 \% de los RF Must tendrá al menos un criterio de aceptación ejecutable y una evidencia de prueba; no se libera si alguno falla. & ISO/IEC/IEEE 29148 \\
RNF-02 & Eficiencia de desempeño & Con 10 000 clientes sintéticos y 20 sesiones concurrentes, el percentil 95 de búsqueda y consulta de membresía será $\leq$ 5 s y la tasa de error $\leq$ 1 \%. & ISO/IEC 25010 \\
RNF-03 & Eficiencia de desempeño & El registro de asistencia confirmada completará el flujo en $\leq$ 30 s para al menos 9 de 10 tareas de una prueba moderada. & ISO/IEC 25010 \\
RNF-04 & Compatibilidad & Las funciones Must superarán el 100 \% de sus pruebas en las dos versiones estables más recientes de Chrome, Edge y Firefox disponibles al corte. & ISO/IEC 25010 \\
RNF-05 & Interacción de usuario & Al menos 9 de 10 tareas críticas de alta, renovación, pago, asistencia y venta se completarán sin ayuda por cada perfil en walkthrough; la aceptación se verificará con registro reproducible por perfil. & ISO/IEC 25010 / Rúbrica \\
RNF-06 & Interacción de usuario & Las pantallas Must cumplirán WCAG 2.2 nivel AA en criterios aplicables mediante auditoría automática y revisión manual, sin errores críticos. & WCAG 2.2 \\
RNF-07 & Fiabilidad & Tras 100 operaciones transaccionales de pago, venta y asistencia con una interrupción simulada, no habrá registros parciales y la recuperación será $\leq$ 10 min. & ISO/IEC 25010 \\
RNF-08 & Fiabilidad & Una copia diaria cifrada permitirá restaurar una base sintética de prueba con RPO $\leq$ 24 h y RTO $\leq$ 4 h en un simulacro documentado. & ISO/IEC 25010 / OWASP \\
RNF-09 & Seguridad & El 100 \% de las pruebas negativas de autorización impedirá acceso a funciones o datos fuera del rol y registrará el intento sin exponer datos personales. & LOPDP / OWASP \\
RNF-10 & Seguridad & Por defecto se almacenarán solo código interno, nombre de uso y contacto opcional; salud, biometría, peso, medidas, dirección e imagen corporal real tendrán cobertura de captura igual a 0 \% en desarrollo y pruebas. & LOPDP / SPDP \\
RNF-11 & Seguridad & Las sesiones expirarán tras 15 min de inactividad; secretos no se almacenarán en texto claro y todo tráfico de un despliegue conectado usará TLS 1.2 o superior. & OWASP / NIST \\
RNF-12 & Mantenibilidad & La suite automatizada cubrirá $\geq$ 70 \% de líneas de los módulos de dominio antes de etiquetar una versión candidata. & ISO/IEC 25010 \\
RNF-13 & Mantenibilidad & Los módulos Clientes, Membresías, Pagos, Asistencia, Inventario, Caja y Rutinas tendrán interfaces explícitas; un cambio de regla en un módulo no exigirá modificar más de un módulo consumidor. & ISO/IEC/IEEE 29148 \\
RNF-14 & Portabilidad & El sistema web y sus pruebas se desplegarán con una única instrucción documentada en Windows 11 y una distribución Linux LTS, sin editar código fuente. & ISO/IEC 25010 \\
RNF-15 & Flexibilidad & Planes, precios, stock mínimo, roles operativos y ventanas de alerta se modificarán por configuración autorizada sin recompilar ni editar código. & ISO/IEC 25010 \\
\bottomrule\end{longtable}
\end{landscape}

\subsection{RNF de explicabilidad derivados y revisados}
Estos requisitos se aplican al \textbf{componente de recomendación propuesto}. Su estado es ESPECIFICADO y no implementado en el MVP actual. La cadena de trazabilidad es evidencia WALK $\rightarrow$ candidato RNF $\rightarrow$ member checking $\rightarrow$ decisión $\rightarrow$ RNF final.
\subsubsection{RNF-16 -- derivado de RNF-EXP-C01}
\begin{tabularx}{\textwidth}{|>{\bfseries}p{4.0cm}|X|}\hline
Enunciado & El componente propuesto de recomendación de rutinas deberá mostrar, junto con cada rutina recomendada, el objetivo o necesidad principal asociada a la rutina y, cuando existan, las restricciones relevantes registradas para el cliente. \\ \hline
Qué se explica & Objetivo/necesidad principal y restricciones relevantes registradas. \\ \hline
A quién & Cliente e instructor autorizado. \\ \hline
Formato & Resumen visible en la vista de la rutina, en lenguaje claro y no técnico. \\ \hline
Momento & Al consultar la recomendación y antes de aceptarla o aplicarla. \\ \hline
Métrica & Porcentaje de recomendaciones que muestran el objetivo/necesidad principal y las restricciones relevantes cuando existan. \\ \hline
Unidad & Porcentaje. \\ \hline
Umbral & 100\% de las recomendaciones generadas en las pruebas deben mostrar el objetivo o necesidad principal y las restricciones registradas cuando existan. \\ \hline
Método de comprobación & Prueba funcional con casos de objetivos y restricciones distintas, verificando la información visible en la rutina. \\ \hline
Responsable & Responsable de calidad con entrenador autorizado. \\ \hline
Frecuencia & En cada versión candidata del componente propuesto. \\ \hline
Fuentes & WALK-TEC-02-F006; WALK-TEC-02-F007; WALK-NTEC-03-F003 \\ \hline
Resultado de member checking & 2 Confirmado; 1 Ajustado \\ \hline
Decisión final & Incorporar\_con\_ajuste \\ \hline
Estado & ESPECIFICADO \\ \hline
\end{tabularx}
\subsubsection{RNF-17 -- derivado de RNF-EXP-C02}
\begin{tabularx}{\textwidth}{|>{\bfseries}p{4.0cm}|X|}\hline
Enunciado & Toda rutina propuesta por el componente inteligente deberá ser revisada por un entrenador autorizado antes de asignarse al cliente. El entrenador podrá aceptar, modificar o rechazar la propuesta; el sistema deberá registrar la decisión, el responsable, la fecha y la versión de la rutina, y mostrar la procedencia profesional externa cuando exista. \\ \hline
Qué se explica & Revisión humana de la rutina y procedencia profesional externa cuando aplique. \\ \hline
A quién & Cliente e instructor autorizado. \\ \hline
Formato & Controles de aceptar/modificar/rechazar para el entrenador y registro auditable de la decisión; etiqueta de estado y procedencia visible para el cliente. \\ \hline
Momento & Antes de que la rutina sea utilizada por el cliente. \\ \hline
Métrica & Porcentaje de propuestas automáticas con decisión humana registrada antes de la asignación y con procedencia profesional visible cuando exista. \\ \hline
Unidad & Porcentaje. \\ \hline
Umbral & 100\% de las propuestas automáticas deben registrar una decisión del entrenador (aceptar, modificar o rechazar) antes de asignarse; toda modificación debe conservar responsable, fecha y versión, y la procedencia profesional debe ser visible cuando exista. \\ \hline
Método de comprobación & Pruebas funcionales con propuestas aceptadas, modificadas y rechazadas; verificar bloqueo de asignación sin revisión y persistencia de decisión, responsable, fecha, versión y procedencia. \\ \hline
Responsable & Entrenador autorizado. \\ \hline
Frecuencia & En cada versión candidata del componente propuesto. \\ \hline
Fuentes & WALK-TEC-02-F008; WALK-TEC-02-F009 \\ \hline
Resultado de member checking & 2 Confirmado; 1 Ajustado \\ \hline
Decisión final & Incorporar\_con\_ajuste \\ \hline
Estado & ESPECIFICADO \\ \hline
\end{tabularx}
\subsubsection{RNF-18 -- derivado de RNF-EXP-C03}
\begin{tabularx}{\textwidth}{|>{\bfseries}p{4.0cm}|X|}\hline
Enunciado & El sistema deberá permitir consultar la rutina vigente y el historial de rutinas registradas, mostrando la fecha y las observaciones asociadas cuando estén disponibles. \\ \hline
Qué se explica & Historial, fecha y observaciones disponibles de las rutinas registradas. \\ \hline
A quién & Cliente e instructor autorizado. \\ \hline
Formato & Historial cronológico de rutinas con fecha y observaciones disponibles. \\ \hline
Momento & Al consultar la rutina vigente o el historial de rutinas. \\ \hline
Métrica & Porcentaje de rutinas registradas consultables con fecha y observaciones cuando estén disponibles. \\ \hline
Unidad & Porcentaje. \\ \hline
Umbral & 100\% de las rutinas registradas durante las pruebas deben ser consultables en el historial con su fecha; las observaciones deben mostrarse cuando existan. \\ \hline
Método de comprobación & Crear y registrar varias rutinas de prueba y verificar su consulta posterior, fecha y observaciones disponibles. \\ \hline
Responsable & Responsable de calidad. \\ \hline
Frecuencia & En cada versión candidata del componente propuesto. \\ \hline
Fuentes & WALK-TEC-01-F003; WALK-TEC-01-F014 \\ \hline
Resultado de member checking & 0 Confirmado; 3 Ajustado \\ \hline
Decisión final & Incorporar\_con\_ajuste \\ \hline
Estado & ESPECIFICADO \\ \hline
\end{tabularx}
\subsubsection{RNF-19 -- derivado de RNF-EXP-C04}
\begin{tabularx}{\textwidth}{|>{\bfseries}p{4.0cm}|X|}\hline
Enunciado & Cuando el sistema sugiera una rutina inicial o realice una asignación automática de una rutina o clase, deberá mostrar de forma breve el criterio principal utilizado para realizar dicha sugerencia o asignación. \\ \hline
Qué se explica & Criterio principal utilizado para una sugerencia u eventual asignación automática. \\ \hline
A quién & Instructor y, cuando corresponda, cliente. \\ \hline
Formato & Texto breve 'Motivo de la recomendación' asociado a la sugerencia. \\ \hline
Momento & En el momento de mostrar la sugerencia y antes de confirmarla. \\ \hline
Métrica & Porcentaje de sugerencias o asignaciones automáticas evaluadas que muestran un criterio principal visible. \\ \hline
Unidad & Porcentaje. \\ \hline
Umbral & 100\% de las sugerencias o asignaciones automáticas evaluadas deben mostrar de forma breve el criterio principal utilizado. \\ \hline
Método de comprobación & Pruebas con perfiles u objetivos distintos, verificando que cada sugerencia o asignación automática muestre el criterio principal utilizado. \\ \hline
Responsable & Responsable de calidad con entrenador autorizado. \\ \hline
Frecuencia & En cada versión candidata del componente propuesto. \\ \hline
Fuentes & WALK-NTEC-01-F008; WALK-TEC-03-F014; WALK-TEC-02-F006 \\ \hline
Resultado de member checking & 0 Confirmado; 3 Ajustado \\ \hline
Decisión final & Incorporar\_con\_ajuste \\ \hline
Estado & ESPECIFICADO \\ \hline
\end{tabularx}


\subsection{RNF complementarios del componente inteligente propuesto}
Los siguientes requisitos proceden de la guía específica y del diseño preventivo, no de observaciones de campo. Sus umbrales son objetivos de verificación futura. Ninguno se presenta como implementado o medido en el MVP.

\subsubsection{RNF-20 -- pertinencia de la recomendación propuesta}
\begin{tabularx}{\textwidth}{|>{\bfseries}p{4.0cm}|X|}\hline
Enunciado & El componente propuesto de recomendación deberá alcanzar una pertinencia validada por entrenador autorizado en al menos 80\% de los casos sintéticos de prueba definidos antes de cada versión candidata. \\ \hline
Métrica / unidad & Casos calificados como pertinentes divididos para casos evaluados / porcentaje. \\ \hline
Umbral & $\geq 80\%$. \\ \hline
Método de verificación & Prueba ciega sobre casos sintéticos predefinidos; un entrenador autorizado registra pertinente/no pertinente. \\ \hline
Responsable / frecuencia & Responsable de calidad con entrenador autorizado / cada versión candidata. \\ \hline
Diseño / prueba & Motor de recomendación propuesto / PV-RNF-20. \\ \hline
Origen / estado & DERIVADO DE NORMA/DISEÑO; PROPUESTO; no existe prueba ejecutada. \\ \hline
\end{tabularx}

\subsubsection{RNF-21 -- equidad de la recomendación propuesta}
\begin{tabularx}{\textwidth}{|>{\bfseries}p{4.0cm}|X|}\hline
Enunciado & En un conjunto de prueba sintético y balanceado, la diferencia máxima entre grupos de edad, condición física declarada y sexo en la tasa de recomendaciones calificadas como no pertinentes por un entrenador autorizado deberá ser menor o igual a 10 puntos porcentuales. \\ \hline
Métrica / unidad & Máxima diferencia absoluta entre tasas por grupo / puntos porcentuales. \\ \hline
Umbral & $\leq 10$ puntos porcentuales. \\ \hline
Método de verificación & Evaluación estratificada con datos sintéticos balanceados y revisión por entrenador. \\ \hline
Responsable / frecuencia & Responsable de calidad con entrenador autorizado / cada versión candidata y ante cambios de datos o reglas. \\ \hline
Diseño / prueba & Módulo de evaluación de equidad propuesto / PV-RNF-21. \\ \hline
Origen / estado & DERIVADO DE NORMA/DISEÑO; PROPUESTO; no existe evidencia de campo ni prueba ejecutada. \\ \hline
\end{tabularx}

\subsubsection{RNF-22 -- monitoreo posterior al despliegue}
\begin{tabularx}{\textwidth}{|>{\bfseries}p{4.0cm}|X|}\hline
Enunciado & Después de un eventual despliegue, el sistema deberá registrar el 100\% de las recomendaciones emitidas, aceptadas, modificadas o rechazadas y producir mensualmente indicadores de aceptación, sustitución humana e incidentes para revisión del responsable del producto. \\ \hline
Métrica / unidad & Eventos con registro completo divididos para eventos emitidos / porcentaje. \\ \hline
Umbral & 100\% de integridad del registro. \\ \hline
Método de verificación & Prueba de integridad del registro y revisión del informe mensual. \\ \hline
Responsable / frecuencia & Responsable del producto / mensual y después de cada cambio de versión. \\ \hline
Diseño / prueba & Registro de monitoreo propuesto / PV-RNF-22. \\ \hline
Origen / estado & DERIVADO DE NORMA/DISEÑO; PROPUESTO; no existe despliegue ni monitoreo ejecutado. \\ \hline
\end{tabularx}

\subsubsection{RNF-23 -- clasificación y control del riesgo}
\begin{tabularx}{\textwidth}{|>{\bfseries}p{4.0cm}|X|}\hline
Enunciado & Antes de habilitar o cambiar el componente propuesto, el responsable del producto deberá clasificar y justificar el 100\% de los riesgos de recomendación en una matriz versionada. El riesgo inicial se clasifica como alto por el posible impacto físico de una recomendación inadecuada; ninguna versión se habilitará sin supervisión humana y aceptación explícita del riesgo residual. \\ \hline
Métrica / unidad & Versiones con evaluación y aprobación de riesgo divididas para versiones candidatas / porcentaje. \\ \hline
Umbral & 100\% antes de habilitación. \\ \hline
Método de verificación & Revisión documental de matriz, controles, justificación y aprobación antes de activar la versión. \\ \hline
Responsable / frecuencia & Responsable del producto y entrenador autorizado / antes de cada habilitación o cambio material. \\ \hline
Diseño / prueba & Matriz de riesgos y control de activación propuestos / PV-RNF-23. \\ \hline
Origen / estado & DERIVADO DE NORMA/DISEÑO; PROPUESTO; no existe evaluación ejecutada. \\ \hline
\end{tabularx}

La operacionalización auditable de RNF-16 a RNF-23 se conserva en \path{04_Trazabilidad/plan_verificacion_rnf_ia.csv}.

\subsection{Cobertura explícita de los seis aspectos del componente inteligente}
La guía de cierre exige seis aspectos del componente inteligente. La cobertura terminal se consolida sin crear identificadores artificiales: algunos aspectos se satisfacen mediante más de un RNF ya vigente. La relación canónica es la siguiente.

\begin{tabularx}{\textwidth}{p{3.0cm}p{2.7cm}X}\toprule
Aspecto exigido & RNF terminal & Criterio verificable y caso de prueba \\ \midrule
Recomendación de rutina & RNF-20 & Pertinencia validada por entrenador $\geq 80\%$ sobre casos sintéticos predefinidos; PV-RNF-20. \\
Explicabilidad & RNF-19 & 100\% de sugerencias automáticas con criterio principal visible; PV-RNF-19. \\
Equidad & RNF-21 & Diferencia máxima de recomendaciones no pertinentes $\leq 10$ puntos porcentuales entre grupos; PV-RNF-21. \\
Supervisión humana & RNF-17 & 100\% de propuestas con decisión del entrenador registrada antes de asignación (aceptar/modificar/rechazar), con responsable, fecha y versión; PV-RNF-17. \\
Monitoreo posterior al despliegue & RNF-22 & 100\% de recomendaciones y decisiones humanas registradas, con resumen mensual; PV-RNF-22. \\
Clasificación de riesgo & RNF-23 & 100\% de versiones con riesgo clasificado y justificado antes de habilitación; PV-RNF-23. \\
\bottomrule
\end{tabularx}

RNF-16 se conserva como requisito complementario de explicación del objetivo y restricciones de la rutina; RNF-18 se conserva como requisito complementario de trazabilidad e historial de rutinas. Para B5 la correspondencia principal de los seis aspectos es uno-a-uno: RNF-20, RNF-19, RNF-21, RNF-17, RNF-22 y RNF-23. El detalle de métricas, unidades, umbrales, responsables, frecuencias, elementos de diseño y casos de prueba está sincronizado en \path{04_Trazabilidad/plan_verificacion_rnf_ia.csv} y \path{04_Trazabilidad/cobertura_6_requisitos_ia.csv}. El componente recomendador permanece propuesto y no se presentan como ejecutadas pruebas que requieren una implementación inexistente.

\subsection{Restricciones de diseño y privacidad}
\begin{tabularx}{\textwidth}{p{2.5cm}X p{3.2cm}}\toprule
ID & Restricción & Fuente \\ \midrule
RD-01 & Excluir datos reales de salud y biometría automatizada; usar minimización, seudonimización y datos sintéticos en diseño y pruebas. & LOPDP / SPDP \\
RD-02 & Aplicar privacidad desde el diseño y por defecto, mínimo privilegio y conservación limitada. & LOPDP / SPDP \\
RD-03 & Prever mecanismos trazables de acceso, rectificación, actualización, eliminación, oposición y portabilidad antes de producción. & LOPDP \\
RD-04 & No incorporar IA o aprendizaje automático al producto mientras no exista necesidad elicitada, evaluación de riesgo, adenda y validación. & Paquete ético / memo de alcance \\
\bottomrule\end{tabularx}

\section{Priorización y alcance del MVP}
La priorización MoSCoW/Kano se conserva del cierre de requisitos. Hay 19 RF Must y 6 RF Should. La lista Must es la referencia para C3. Las fuentes históricas presentan una contradicción entre 12/19 y 16/19 RF Must cubiertos. Esta ERS no resuelve esa diferencia por inferencia: la cobertura MVP queda \textbf{sujeta a verificación terminal C3} contra el código y las pruebas antes de marcar cada RF como IMPLEMENTADO.

\small
\setlength{\tabcolsep}{4pt}
\begin{longtable}{p{1.8cm}p{1.45cm}p{1.65cm}p{1.0cm}p{7.7cm}}
\toprule RF & MoSCoW & Kano & WSJF & Objetivo \\ \midrule
\endfirsthead
\toprule RF & MoSCoW & Kano & WSJF & Objetivo \\ \midrule
\endhead
RF-01 & Must & Básico & 9.67 & El sistema deberá permitir iniciar y cerrar sesión mediante credenciales individuales no biométricas. \\
RF-05 & Must & Desempeño & 7.00 & El sistema deberá buscar clientes por código, nombre normalizado o contacto y mostrar membresía, último pago, asistencia y novedades según permisos. \\
RF-09 & Must & Básico & 7.00 & El sistema deberá mostrar estado, plan, fecha de inicio y vencimiento de la membresía. \\
RF-08 & Must & Básico & 6.75 & El sistema deberá activar o renovar una membresía calculando inicio y vencimiento desde un plan vigente y un pago confirmado. \\
RF-11 & Must & Básico & 6.75 & El sistema deberá registrar monto, concepto, medio, referencia y fecha de un pago, y generar un comprobante interno. \\
RF-04 & Must & Desempeño & 6.00 & El sistema deberá registrar un cliente con código interno, nombre de uso y un medio de contacto, sin salud, biometría, dirección ni imagen corporal. \\
RF-13 & Must & Básico & 6.00 & El sistema deberá registrar una asistencia no biométrica en una acción tras validar identidad operativa y vigencia, o documentar una excepción autorizada. \\
RF-14 & Must & Desempeño & 6.00 & El sistema deberá filtrar asistencias por cliente, fecha y turno, y mostrar conteos operativos. \\
RF-17 & Must & Básico & 6.00 & El sistema deberá registrar una venta con ítems y medio de pago, y descontar existencias en la misma transacción. \\
RF-10 & Must & Desempeño & 5.25 & El sistema deberá listar diariamente membresías que vencen en los próximos tres días y las vencidas, sin enviar mensajes externos automáticamente. \\
RF-02 & Must & Básico & 5.20 & El sistema deberá autorizar cada operación según el rol Administrador, Recepción o Instructor y denegar por defecto las no asignadas. \\
RF-15 & Must & Desempeño & 5.00 & El sistema deberá registrar productos con código, nombre, precio, unidad, stock mínimo y estado. \\
RF-19 & Must & Desempeño & 4.80 & El sistema deberá calcular pagos y ventas por medio, comparar valores declarados y registrar diferencias y entrega de turno. \\
RF-06 & Must & Desempeño & 4.75 & El sistema deberá actualizar o reactivar un registro existente y advertir coincidencias antes de crear otro cliente. \\
RF-20 & Must & Desempeño & 4.75 & El sistema deberá crear una novedad con categoría, detalle mínimo, responsable, vencimiento y estados pendiente, en proceso, resuelta o cancelada. \\
RF-23 & Must & Desempeño & 4.75 & El sistema deberá crear una nueva versión al modificar ejercicios o parámetros y conservar las versiones anteriores en solo lectura. \\
RF-16 & Must & Desempeño & 4.60 & El sistema deberá registrar entradas y ajustes con cantidad, motivo, responsable y saldo resultante. \\
RF-07 & Must & Desempeño & 4.40 & El sistema deberá administrar nombre, duración, precio, vigencia y estado de planes y promociones. \\
RF-22 & Must & Básico & 4.20 & El sistema deberá crear una rutina con objetivo operativo, ejercicios, series, repeticiones, descanso y días, y asignarla a un cliente sintético o autorizado. \\
RF-03 & Should & Desempeño &  & El sistema deberá registrar autor, fecha, acción y valores anterior/nuevo para pagos, precios, inventario, membresías y permisos. \\
RF-12 & Should & Desempeño &  & El sistema deberá permitir confirmar, rechazar o corregir un pago pendiente conservando motivo y auditoría. \\
RF-18 & Should & Desempeño &  & El sistema deberá listar productos en o bajo stock mínimo y resumir entradas, salidas y ventas por período. \\
RF-21 & Should & Desempeño &  & El sistema deberá registrar turnos, capacidad, disponibilidad y novedades operativas del instructor, sin datos de salud. \\
RF-24 & Should & Desempeño &  & El sistema deberá registrar fecha, cumplimiento de ejercicios, repeticiones y observación operativa, excluyendo diagnósticos, lesiones, peso y medidas reales. \\
RF-25 & Should & Desempeño &  & El sistema deberá mostrar indicadores filtrables de membresías, pagos, asistencias, ventas, caja, inventario y novedades, con exportación CSV. \\
\bottomrule
\caption{Priorización terminal de los 25 RF.}
\end{longtable}
\setlength{\tabcolsep}{6pt}
\normalsize

\section{Historias de usuario y criterios de aceptación Must}
Las 19 historias Must mantienen correspondencia uno a uno con CU-01..CU-19 y CA-01..CA-19. Para conservar un artefacto ejecutable, el criterio se expresa como escenario observable y puede materializarse en Gherkin en el paquete de pruebas.
\subsubsection{HU-01 / CU-01 / RF-01}
\textbf{Connextra.} Como Personal autorizado, quiero garantizar tratamiento lícito y minimización de datos, para completar la operación de forma trazable y consistente.\\
\textbf{Gherkin / CA-01.} Dadas credenciales de demostración válidas, cuando el usuario inicia sesión, entonces se crea una sesión del rol correspondiente; con credenciales inválidas se deniega el acceso sin revelar qué campo falló.
\subsubsection{HU-02 / CU-02 / RF-02}
\textbf{Connextra.} Como Administrador, quiero controlar acceso por roles y permisos, para completar la operación de forma trazable y consistente.\\
\textbf{Gherkin / CA-02.} Dado un usuario autenticado, cuando intenta una operación fuera de su rol, entonces el sistema la deniega; cuando la operación está autorizada, permite continuar.
\subsubsection{HU-03 / CU-03 / RF-04}
\textbf{Connextra.} Como Recepción, quiero registrar datos mínimos de cliente, para completar la operación de forma trazable y consistente.\\
\textbf{Gherkin / CA-03.} Dado un operador autorizado, cuando registra código interno, nombre de uso y contacto permitido, entonces se crea un cliente único sin solicitar salud, biometría, dirección ni imagen corporal.
\subsubsection{HU-04 / CU-04 / RF-05}
\textbf{Connextra.} Como Recepción, quiero consultar ficha de cliente, para completar la operación de forma trazable y consistente.\\
\textbf{Gherkin / CA-04.} Dado un criterio de búsqueda válido, cuando se consulta un cliente, entonces se muestra su ficha operativa permitida sin exponer información fuera del rol.
\subsubsection{HU-05 / CU-05 / RF-06}
\textbf{Connextra.} Como Recepción, quiero actualizar y reactivar cliente, para completar la operación de forma trazable y consistente.\\
\textbf{Gherkin / CA-05.} Dado un cliente existente, cuando se actualiza o reactiva, entonces se conserva el mismo identificador y se advierten coincidencias antes de crear otro registro.
\subsubsection{HU-06 / CU-06 / RF-07}
\textbf{Connextra.} Como Administrador, quiero configurar catálogo de planes, para completar la operación de forma trazable y consistente.\\
\textbf{Gherkin / CA-06.} Dado un administrador, cuando crea o modifica un plan, entonces quedan registrados nombre, duración, precio, vigencia y estado y el catálogo refleja el cambio.
\subsubsection{HU-07 / CU-07 / RF-08}
\textbf{Connextra.} Como Recepción, quiero activar y renovar membresía, para completar la operación de forma trazable y consistente.\\
\textbf{Gherkin / CA-07.} Dado un plan vigente y un pago confirmado, cuando se activa o renueva una membresía, entonces el inicio y vencimiento se calculan y persisten de forma consistente.
\subsubsection{HU-08 / CU-08 / RF-09}
\textbf{Connextra.} Como Personal autorizado, quiero consultar vigencia, para completar la operación de forma trazable y consistente.\\
\textbf{Gherkin / CA-08.} Dado un cliente localizado, cuando se consulta su membresía, entonces se muestran estado, plan, inicio y vencimiento.
\subsubsection{HU-09 / CU-09 / RF-10}
\textbf{Connextra.} Como Recepción, quiero alertar vencimientos, para completar la operación de forma trazable y consistente.\\
\textbf{Gherkin / CA-09.} Dada la fecha de corte, cuando se consulta la alerta, entonces se listan las membresías vencidas y las que vencen dentro de tres días sin enviar mensajes externos automáticamente.
\subsubsection{HU-10 / CU-10 / RF-11}
\textbf{Connextra.} Como Recepción, quiero registrar pago y comprobante, para completar la operación de forma trazable y consistente.\\
\textbf{Gherkin / CA-10.} Dado un pago válido, cuando se registra, entonces se conservan monto, concepto, medio, referencia y fecha y se genera un comprobante interno.
\subsubsection{HU-11 / CU-11 / RF-13}
\textbf{Connextra.} Como Recepción, quiero registrar asistencia, para completar la operación de forma trazable y consistente.\\
\textbf{Gherkin / CA-11.} Dado un cliente con vigencia válida, cuando se registra asistencia, entonces se crea un único evento no biométrico; si requiere excepción, se documenta de forma autorizada.
\subsubsection{HU-12 / CU-12 / RF-14}
\textbf{Connextra.} Como Personal autorizado, quiero consultar asistencia, para completar la operación de forma trazable y consistente.\\
\textbf{Gherkin / CA-12.} Dado un usuario autorizado, cuando filtra asistencias por cliente, fecha o turno, entonces se muestran el detalle y los conteos correspondientes.
\subsubsection{HU-13 / CU-13 / RF-15}
\textbf{Connextra.} Como Administrador, quiero administrar productos, para completar la operación de forma trazable y consistente.\\
\textbf{Gherkin / CA-13.} Dado un administrador, cuando registra un producto, entonces se almacenan código, nombre, precio, unidad, stock mínimo y estado.
\subsubsection{HU-14 / CU-14 / RF-16}
\textbf{Connextra.} Como Personal autorizado, quiero registrar movimientos de stock, para completar la operación de forma trazable y consistente.\\
\textbf{Gherkin / CA-14.} Dado un producto existente, cuando se registra una entrada o ajuste, entonces se guardan cantidad, motivo, responsable y saldo resultante.
\subsubsection{HU-15 / CU-15 / RF-17}
\textbf{Connextra.} Como Recepción, quiero registrar venta, para completar la operación de forma trazable y consistente.\\
\textbf{Gherkin / CA-15.} Dados ítems con stock suficiente, cuando se registra una venta, entonces la venta y el descuento de existencias ocurren en la misma operación demostrativa.
\subsubsection{HU-16 / CU-16 / RF-19}
\textbf{Connextra.} Como Recepción, quiero cerrar turno y conciliar caja, para completar la operación de forma trazable y consistente.\\
\textbf{Gherkin / CA-16.} Dado un turno abierto, cuando se ejecuta el cierre, entonces se calculan pagos y ventas por medio, se compara el valor declarado y se registra la diferencia.
\subsubsection{HU-17 / CU-17 / RF-20}
\textbf{Connextra.} Como Personal autorizado, quiero gestionar novedades internas, para completar la operación de forma trazable y consistente.\\
\textbf{Gherkin / CA-17.} Dado un usuario autorizado, cuando crea una novedad, entonces se guardan categoría, detalle, responsable, vencimiento y estado y se conserva el historial de estado.
\subsubsection{HU-18 / CU-18 / RF-22}
\textbf{Connextra.} Como Instructor, quiero administrar rutinas, para completar la operación de forma trazable y consistente.\\
\textbf{Gherkin / CA-18.} Dado un instructor autorizado, cuando crea y asigna una rutina, entonces se registran objetivo, ejercicios, series, repeticiones, descanso y días sin datos de salud reales.
\subsubsection{HU-19 / CU-19 / RF-23}
\textbf{Connextra.} Como Instructor, quiero registrar seguimiento, para completar la operación de forma trazable y consistente.\\
\textbf{Gherkin / CA-19.} Dada una rutina existente, cuando se modifica un ejercicio o parámetro, entonces se crea una nueva versión y las versiones anteriores permanecen en solo lectura.


\section{Especificación textual de los 19 casos de uso Must}
Esta sección complementa los diagramas UML con la especificación de comportamiento exigida para el cierre del examen suspenso. Cada caso de uso se deriva del RF, la HU, el criterio de aceptación y la evidencia ya versionados; no introduce funcionalidad nueva. Los identificadores de flujo \texttt{CU-xx-FP}, \texttt{CU-xx-FA-01} y \texttt{CU-xx-EX-01} se trazan de forma explícita en \path{04_Trazabilidad/matriz_trazabilidad.csv}.

\begin{importantbox}
\textbf{Convención de trazabilidad de flujos.} \texttt{FP} identifica el flujo principal, \texttt{FA} un flujo alternativo válido del negocio y \texttt{EX} una excepción que impide o desvía la operación. Cada \texttt{FA} declara de forma explícita el paso del \texttt{FP} desde el cual se abre y el punto al que retorna o si termina el caso de uso. Las condiciones y respuestas aquí descritas refinan el comportamiento de los requisitos Must existentes; no crean RF adicionales.
\end{importantbox}

\subsubsection{CU-01 -- Autenticar usuario}
\begin{tabularx}{\textwidth}{|>{\bfseries}p{3.2cm}|X|}\hline
Actor principal & Personal autorizado (Administrador, Recepción o Instructor) \\ \hline
Requisito relacionado & \texttt{RF-01} \\ \hline
Historia / criterio & \texttt{HU-01} / \texttt{CA-01} \\ \hline
Componente & Seguridad \\ \hline
Evidencia de origen & {\ttfamily\seqsplit{EV-ENTR02-07}} \\ \hline
Mockups relacionados & {\ttfamily\seqsplit{MU-01}} \\ \hline
Disparador & El usuario solicita acceso al sistema mediante sus credenciales individuales. \\ \hline
\end{tabularx}

\textbf{Precondiciones.}
\begin{itemize}
\item El usuario puede acceder a la interfaz de autenticación.
\item La interfaz de autenticación y el servicio de seguridad están disponibles.
\end{itemize}

\textbf{Flujo principal \texttt{CU-01-FP}.}
\begin{enumerate}
\item El usuario abre la pantalla de acceso.
\item El usuario ingresa sus credenciales individuales no biométricas.
\item El sistema valida que los campos requeridos estén presentes.
\item El servicio de seguridad contrasta las credenciales con el registro autorizado.
\item El sistema identifica el rol asociado a la cuenta válida.
\item El sistema crea la sesión correspondiente al rol.
\item El sistema registra el evento de autenticación requerido y muestra la vista inicial sin exponer datos excesivos.
\end{enumerate}

\textbf{Flujo alternativo \texttt{CU-01-FA-01}.} \textit{Condición de disparo:} En el paso 3 del flujo principal falta uno de los campos requeridos para intentar la autenticación. \textit{Apertura y retorno:} Se abre desde el paso 3; después de completar el campo faltante, retorna al paso 3 para repetir la validación.
\begin{enumerate}[label=\alph*.]
\item El sistema identifica el campo requerido que falta sin validar todavía las credenciales.
\item El usuario completa el campo requerido.
\end{enumerate}

\textbf{Excepción \texttt{CU-01-EX-01}.} \textit{Condición de disparo:} Las credenciales son inválidas o la cuenta no está habilitada.
\begin{enumerate}[label=\alph*.]
\item El sistema deniega el acceso con un mensaje genérico.
\item El sistema no revela cuál campo falló ni crea una sesión.
\end{enumerate}

\textbf{Poscondiciones.}
\begin{itemize}
\item Con acceso exitoso, existe una sesión activa asociada al rol autorizado.
\item Con rechazo de la autenticación, no se crea una sesión activa.
\end{itemize}

\subsubsection{CU-02 -- Aplicar permisos por rol}
\begin{tabularx}{\textwidth}{|>{\bfseries}p{3.2cm}|X|}\hline
Actor principal & Administrador \\ \hline
Requisito relacionado & \texttt{RF-02} \\ \hline
Historia / criterio & \texttt{HU-02} / \texttt{CA-02} \\ \hline
Componente & Seguridad \\ \hline
Evidencia de origen & {\ttfamily\seqsplit{EV-ENTR02-07;EV-ENTR03-06;EV-ENTR08-04}} \\ \hline
Mockups relacionados & {\ttfamily\seqsplit{MU-22;MU-32}} \\ \hline
Disparador & Un usuario autenticado solicita ejecutar una operación protegida. \\ \hline
\end{tabularx}

\textbf{Precondiciones.}
\begin{itemize}
\item Existe una sesión autenticada.
\item El usuario tiene asignado uno de los roles Administrador, Recepción o Instructor.
\end{itemize}

\textbf{Flujo principal \texttt{CU-02-FP}.}
\begin{enumerate}
\item El usuario solicita una operación desde la interfaz.
\item El sistema identifica el rol de la sesión activa.
\item El servicio de seguridad consulta las reglas y permisos asociados al rol.
\item El sistema comprueba que la operación está autorizada.
\item El sistema habilita la operación y continúa con el caso de uso solicitado.
\item El sistema registra el evento de autorización cuando corresponde.
\end{enumerate}

\textbf{Flujo alternativo \texttt{CU-02-FA-01}.} \textit{Condición de disparo:} El rol permite consulta pero no modificación sobre el recurso solicitado. \textit{Apertura y retorno:} Se abre desde el paso 4 del flujo principal; el caso de uso termina al presentar el recurso en modo de consulta, sin ejecutar el paso 5 como operación de modificación.
\begin{enumerate}[label=\alph*.]
\item El sistema presenta el recurso en modo de consulta.
\item Las acciones de modificación permanecen deshabilitadas para ese rol.
\end{enumerate}

\textbf{Excepción \texttt{CU-02-EX-01}.} \textit{Condición de disparo:} La operación no está asignada al rol o no existe una regla explícita que la autorice.
\begin{enumerate}[label=\alph*.]
\item El sistema aplica denegación por defecto.
\item La operación no se ejecuta y se informa la denegación sin exponer reglas internas.
\end{enumerate}

\textbf{Poscondiciones.}
\begin{itemize}
\item Solo se ejecutan operaciones compatibles con el rol autenticado.
\item Una operación denegada no modifica información del sistema.
\end{itemize}

\subsubsection{CU-03 -- Registrar cliente mínimo}
\begin{tabularx}{\textwidth}{|>{\bfseries}p{3.2cm}|X|}\hline
Actor principal & Recepción \\ \hline
Requisito relacionado & \texttt{RF-04} \\ \hline
Historia / criterio & \texttt{HU-03} / \texttt{CA-03} \\ \hline
Componente & Clientes \\ \hline
Evidencia de origen & {\ttfamily\seqsplit{EV-ENTR01-02;EV-ENTR02-02;EV-ENTR03-02;EV-ENTR08-01}} \\ \hline
Mockups relacionados & {\ttfamily\seqsplit{MU-12;MU-37}} \\ \hline
Disparador & Recepción solicita crear un cliente. \\ \hline
\end{tabularx}

\textbf{Precondiciones.}
\begin{itemize}
\item El operador está autenticado y autorizado para gestionar clientes.
\item El formulario de registro mínimo está disponible para el operador.
\end{itemize}

\textbf{Flujo principal \texttt{CU-03-FP}.}
\begin{enumerate}
\item Recepción abre el formulario de registro mínimo.
\item Ingresa código interno, nombre de uso y medio de contacto.
\item El sistema valida formato y obligatoriedad de los datos mínimos.
\item El sistema comprueba que el código interno no esté registrado.
\item El sistema crea el registro del cliente sin solicitar datos de salud, biometría, dirección ni imagen corporal.
\item El sistema confirma la creación y deja disponible el código del cliente.
\end{enumerate}

\textbf{Flujo alternativo \texttt{CU-03-FA-01}.} \textit{Condición de disparo:} En el paso 3 del flujo principal los datos mínimos están presentes, pero el medio de contacto no cumple el formato permitido. \textit{Apertura y retorno:} Se abre desde el paso 3; una vez corregido el contacto, retorna al paso 3 para repetir la validación de formato y obligatoriedad.
\begin{enumerate}[label=\alph*.]
\item El sistema identifica el medio de contacto que requiere corrección y conserva los demás datos ya ingresados.
\item Recepción corrige el medio de contacto sin añadir categorías de datos excluidas por el alcance.
\end{enumerate}

\textbf{Excepción \texttt{CU-03-EX-01}.} \textit{Condición de disparo:} Falta un dato mínimo obligatorio o el código interno ya existe.
\begin{enumerate}[label=\alph*.]
\item El sistema no crea un nuevo cliente.
\item Se informa la validación necesaria sin incorporar datos adicionales al expediente.
\end{enumerate}

\textbf{Poscondiciones.}
\begin{itemize}
\item Existe un cliente único con el conjunto mínimo permitido de datos.
\item No se incorporan categorías de datos excluidas por el alcance.
\end{itemize}

\subsubsection{CU-04 -- Buscar y consultar cliente}
\begin{tabularx}{\textwidth}{|>{\bfseries}p{3.2cm}|X|}\hline
Actor principal & Recepción \\ \hline
Requisito relacionado & \texttt{RF-05} \\ \hline
Historia / criterio & \texttt{HU-04} / \texttt{CA-04} \\ \hline
Componente & Clientes \\ \hline
Evidencia de origen & {\ttfamily\seqsplit{EV-ENTR02-03;EV-ENTR02-08;EV-ENTR08-05}} \\ \hline
Mockups relacionados & {\ttfamily\seqsplit{MU-11;MU-36}} \\ \hline
Disparador & Recepción solicita localizar y consultar una ficha de cliente. \\ \hline
\end{tabularx}

\textbf{Precondiciones.}
\begin{itemize}
\item El usuario está autenticado con permiso de consulta.
\item Existe un criterio de búsqueda por código, nombre normalizado o contacto.
\end{itemize}

\textbf{Flujo principal \texttt{CU-04-FP}.}
\begin{enumerate}
\item El usuario introduce un criterio de búsqueda.
\item El sistema normaliza y valida el criterio.
\item El servicio de clientes consulta coincidencias autorizadas.
\item El sistema presenta el cliente localizado.
\item El usuario selecciona la ficha cuando corresponde.
\item El sistema muestra membresía, último pago, asistencia y novedades permitidas por el rol.
\end{enumerate}

\textbf{Flujo alternativo \texttt{CU-04-FA-01}.} \textit{Condición de disparo:} La búsqueda devuelve más de una coincidencia válida. \textit{Apertura y retorno:} Se abre desde el paso 3 del flujo principal; después de que el usuario selecciona una coincidencia, retorna al paso 6 para mostrar la ficha autorizada.
\begin{enumerate}[label=\alph*.]
\item El sistema muestra una lista acotada de coincidencias permitidas.
\item El usuario selecciona el cliente y el sistema continúa con la consulta de su ficha.
\end{enumerate}

\textbf{Excepción \texttt{CU-04-EX-01}.} \textit{Condición de disparo:} No existe coincidencia o el rol no puede consultar la información solicitada.
\begin{enumerate}[label=\alph*.]
\item El sistema informa que no hay resultados disponibles para la consulta autorizada.
\item No se revelan datos de clientes fuera del alcance del rol.
\end{enumerate}

\textbf{Poscondiciones.}
\begin{itemize}
\item La consulta no modifica el registro del cliente.
\item Solo se visualiza información operativa permitida.
\end{itemize}

\subsubsection{CU-05 -- Actualizar o reactivar cliente}
\begin{tabularx}{\textwidth}{|>{\bfseries}p{3.2cm}|X|}\hline
Actor principal & Recepción \\ \hline
Requisito relacionado & \texttt{RF-06} \\ \hline
Historia / criterio & \texttt{HU-05} / \texttt{CA-05} \\ \hline
Componente & Clientes \\ \hline
Evidencia de origen & {\ttfamily\seqsplit{EV-ENTR02-02;EV-ENTR03-02;EV-ENTR04-01;EV-ENTR05-02}} \\ \hline
Mockups relacionados & {\ttfamily\seqsplit{MU-11;MU-36}} \\ \hline
Disparador & Recepción solicita modificar o reactivar un cliente existente. \\ \hline
\end{tabularx}

\textbf{Precondiciones.}
\begin{itemize}
\item El usuario está autorizado para gestionar clientes.
\item El cliente objetivo existe y puede ser localizado.
\end{itemize}

\textbf{Flujo principal \texttt{CU-05-FP}.}
\begin{enumerate}
\item El usuario busca y selecciona el cliente existente.
\item El sistema carga el registro con su identificador actual.
\item El usuario modifica los datos permitidos o solicita reactivación.
\item El sistema valida los cambios y comprueba posibles coincidencias.
\item El sistema actualiza el mismo registro sin crear un nuevo identificador.
\item El sistema confirma la actualización o reactivación.
\end{enumerate}

\textbf{Flujo alternativo \texttt{CU-05-FA-01}.} \textit{Condición de disparo:} En el paso 3 del flujo principal el cliente seleccionado se encuentra inactivo y el usuario solicita reactivarlo. \textit{Apertura y retorno:} Se abre desde el paso 3; después de validar la reactivación sobre el mismo identificador, retorna al paso 6 para confirmar el resultado.
\begin{enumerate}[label=\alph*.]
\item El sistema conserva el identificador existente y prepara el cambio de estado del cliente.
\item El sistema valida que la reactivación no genere una duplicidad y reactiva el mismo registro.
\end{enumerate}

\textbf{Excepción \texttt{CU-05-EX-01}.} \textit{Condición de disparo:} Los cambios producirían una duplicidad o contienen datos no válidos.
\begin{enumerate}[label=\alph*.]
\item El sistema rechaza la actualización conflictiva.
\item El registro original permanece sin cambios.
\end{enumerate}

\textbf{Poscondiciones.}
\begin{itemize}
\item El cliente conserva su identificador original.
\item La operación no crea duplicados.
\end{itemize}

\subsubsection{CU-06 -- Configurar planes y promociones}
\begin{tabularx}{\textwidth}{|>{\bfseries}p{3.2cm}|X|}\hline
Actor principal & Administrador \\ \hline
Requisito relacionado & \texttt{RF-07} \\ \hline
Historia / criterio & \texttt{HU-06} / \texttt{CA-06} \\ \hline
Componente & Membresías \\ \hline
Evidencia de origen & {\ttfamily\seqsplit{EV-ENTR02-01;EV-ENTR04-01}} \\ \hline
Mockups relacionados & {\ttfamily\seqsplit{MU-68}} \\ \hline
Disparador & El administrador solicita crear o modificar un plan o promoción. \\ \hline
\end{tabularx}

\textbf{Precondiciones.}
\begin{itemize}
\item El administrador mantiene una sesión válida.
\item La operación de configuración de catálogo está autorizada.
\end{itemize}

\textbf{Flujo principal \texttt{CU-06-FP}.}
\begin{enumerate}
\item El administrador abre el catálogo de planes y promociones.
\item Ingresa o modifica nombre, duración, precio, vigencia y estado.
\item El sistema valida los valores ingresados.
\item El sistema guarda la configuración del plan o promoción.
\item El catálogo refleja inmediatamente el estado vigente del registro.
\item El sistema confirma la operación y registra el evento requerido.
\end{enumerate}

\textbf{Flujo alternativo \texttt{CU-06-FA-01}.} \textit{Condición de disparo:} El administrador decide desactivar un plan o promoción existente sin eliminarlo. \textit{Apertura y retorno:} Se abre desde el paso 2 del flujo principal; después de cambiar el estado a inactivo, retorna al paso 5 para reflejar el estado actualizado en el catálogo.
\begin{enumerate}[label=\alph*.]
\item El sistema cambia el estado del registro a inactivo.
\item El registro permanece disponible para trazabilidad histórica y deja de ofrecerse como opción vigente.
\end{enumerate}

\textbf{Excepción \texttt{CU-06-EX-01}.} \textit{Condición de disparo:} La duración, el precio o el periodo de vigencia son inválidos.
\begin{enumerate}[label=\alph*.]
\item El sistema no guarda la configuración.
\item Se señalan los campos que requieren corrección sin alterar el catálogo vigente.
\end{enumerate}

\textbf{Poscondiciones.}
\begin{itemize}
\item El catálogo conserva una configuración consistente de planes y promociones.
\item Los registros inactivos no se presentan como planes vigentes.
\end{itemize}

\subsubsection{CU-07 -- Activar o renovar membresía}
\begin{tabularx}{\textwidth}{|>{\bfseries}p{3.2cm}|X|}\hline
Actor principal & Recepción \\ \hline
Requisito relacionado & \texttt{RF-08} \\ \hline
Historia / criterio & \texttt{HU-07} / \texttt{CA-07} \\ \hline
Componente & Membresías \\ \hline
Evidencia de origen & {\ttfamily\seqsplit{EV-ENTR01-02;EV-ENTR02-04;EV-ENTR03-01;EV-ENTR04-03}} \\ \hline
Mockups relacionados & {\ttfamily\seqsplit{MU-14;MU-39}} \\ \hline
Disparador & Recepción solicita activar o renovar la membresía de un cliente. \\ \hline
\end{tabularx}

\textbf{Precondiciones.}
\begin{itemize}
\item Recepción está autenticada y autorizada para gestionar membresías.
\item El cliente existe y puede ser localizado para iniciar la operación.
\end{itemize}

\textbf{Flujo principal \texttt{CU-07-FP}.}
\begin{enumerate}
\item Recepción localiza al cliente.
\item Selecciona el plan vigente aplicable.
\item El sistema verifica la confirmación del pago.
\item El sistema determina la fecha de inicio correspondiente.
\item El sistema calcula la fecha de vencimiento según la duración del plan.
\item El sistema registra la membresía y sus fechas.
\item El sistema confirma la activación.
\end{enumerate}

\textbf{Flujo alternativo \texttt{CU-07-FA-01}.} \textit{Condición de disparo:} El cliente ya posee una membresía y la operación corresponde a renovación. \textit{Apertura y retorno:} Se abre desde el paso 4 del flujo principal; después de determinar el inicio de la renovación conservando el historial, retorna al paso 5 para calcular el nuevo vencimiento.
\begin{enumerate}[label=\alph*.]
\item El sistema identifica la membresía vigente o previa que será renovada y conserva su historial.
\item El sistema determina la fecha de inicio de la renovación sin sobrescribir los periodos anteriores.
\end{enumerate}

\textbf{Excepción \texttt{CU-07-EX-01}.} \textit{Condición de disparo:} El plan no está vigente o no existe pago confirmado.
\begin{enumerate}[label=\alph*.]
\item El sistema impide activar o renovar la membresía.
\item No se modifican las fechas de vigencia existentes.
\end{enumerate}

\textbf{Poscondiciones.}
\begin{itemize}
\item La activación o renovación queda asociada al cliente, plan y pago confirmado.
\item Las fechas de inicio y vencimiento quedan persistidas de forma consistente.
\end{itemize}

\subsubsection{CU-08 -- Consultar vigencia}
\begin{tabularx}{\textwidth}{|>{\bfseries}p{3.2cm}|X|}\hline
Actor principal & Personal autorizado \\ \hline
Requisito relacionado & \texttt{RF-09} \\ \hline
Historia / criterio & \texttt{HU-08} / \texttt{CA-08} \\ \hline
Componente & Membresías \\ \hline
Evidencia de origen & {\ttfamily\seqsplit{EV-ENTR01-03;EV-ENTR02-03;EV-ENTR03-03;EV-ENTR08-02}} \\ \hline
Mockups relacionados & {\ttfamily\seqsplit{MU-13;MU-38;MU-63(F)}} \\ \hline
Disparador & Un usuario autorizado solicita consultar la vigencia de una membresía. \\ \hline
\end{tabularx}

\textbf{Precondiciones.}
\begin{itemize}
\item El usuario mantiene una sesión autorizada para consultar membresías.
\item La interfaz de búsqueda de clientes está disponible.
\end{itemize}

\textbf{Flujo principal \texttt{CU-08-FP}.}
\begin{enumerate}
\item El usuario busca al cliente.
\item El sistema recupera la membresía asociada.
\item El sistema determina el estado de la membresía con base en sus fechas.
\item El sistema muestra estado, plan, fecha de inicio y fecha de vencimiento.
\item La consulta queda disponible sin modificar el registro.
\end{enumerate}

\textbf{Flujo alternativo \texttt{CU-08-FA-01}.} \textit{Condición de disparo:} El cliente no posee una membresía activa en la fecha de consulta. \textit{Apertura y retorno:} Se abre desde el paso 3 del flujo principal; el caso de uso termina después de presentar el estado no vigente, sin ejecutar los pasos destinados a mostrar una membresía activa.
\begin{enumerate}[label=\alph*.]
\item El sistema muestra el estado no vigente correspondiente.
\item No se presenta una membresía inexistente como activa.
\end{enumerate}

\textbf{Excepción \texttt{CU-08-EX-01}.} \textit{Condición de disparo:} El cliente no existe o el usuario carece de permiso de consulta.
\begin{enumerate}[label=\alph*.]
\item El sistema no expone información de membresía.
\item Se informa que la consulta no puede completarse.
\end{enumerate}

\textbf{Poscondiciones.}
\begin{itemize}
\item La operación es de solo lectura.
\item La información presentada corresponde al alcance del rol.
\end{itemize}

\subsubsection{CU-09 -- Alertar vencimientos}
\begin{tabularx}{\textwidth}{|>{\bfseries}p{3.2cm}|X|}\hline
Actor principal & Recepción \\ \hline
Requisito relacionado & \texttt{RF-10} \\ \hline
Historia / criterio & \texttt{HU-09} / \texttt{CA-09} \\ \hline
Componente & Membresías \\ \hline
Evidencia de origen & {\ttfamily\seqsplit{EV-ENTR01-03;EV-ENTR05-02;EV-ENTR07-02}} \\ \hline
Mockups relacionados & {\ttfamily\seqsplit{MU-10;MU-13;MU-38}} \\ \hline
Disparador & Recepción consulta las membresías que requieren atención por vencimiento. \\ \hline
\end{tabularx}

\textbf{Precondiciones.}
\begin{itemize}
\item El usuario está autenticado y autorizado para consultar vencimientos.
\item La vista de vencimientos está disponible para iniciar la consulta.
\end{itemize}

\textbf{Flujo principal \texttt{CU-09-FP}.}
\begin{enumerate}
\item El usuario solicita la vista de vencimientos.
\item El sistema toma la fecha de corte.
\item El sistema consulta membresías vencidas y próximas a vencer.
\item El sistema identifica las que vencen dentro de tres días.
\item El sistema ordena y presenta una lista priorizada.
\item El sistema no envía mensajes externos automáticamente.
\end{enumerate}

\textbf{Flujo alternativo \texttt{CU-09-FA-01}.} \textit{Condición de disparo:} No existen membresías vencidas ni próximas a vencer en la ventana definida. \textit{Apertura y retorno:} Se abre desde el paso 3 del flujo principal; el caso de uso termina después de presentar el resultado vacío con la fecha de corte aplicada.
\begin{enumerate}[label=\alph*.]
\item El sistema presenta una lista vacía con la fecha de corte aplicada.
\item No se genera comunicación externa.
\end{enumerate}

\textbf{Excepción \texttt{CU-09-EX-01}.} \textit{Condición de disparo:} La fecha de corte o el filtro de consulta es inválido.
\begin{enumerate}[label=\alph*.]
\item El sistema no ejecuta la búsqueda.
\item Se solicita corregir el criterio antes de volver a consultar.
\end{enumerate}

\textbf{Poscondiciones.}
\begin{itemize}
\item Se obtiene una lista de atención basada en vigencia.
\item La consulta no modifica membresías ni dispara mensajería externa.
\end{itemize}

\subsubsection{CU-10 -- Registrar pago y comprobante}
\begin{tabularx}{\textwidth}{|>{\bfseries}p{3.2cm}|X|}\hline
Actor principal & Recepción \\ \hline
Requisito relacionado & \texttt{RF-11} \\ \hline
Historia / criterio & \texttt{HU-10} / \texttt{CA-10} \\ \hline
Componente & Pagos \\ \hline
Evidencia de origen & {\ttfamily\seqsplit{EV-ENTR01-02;EV-ENTR02-04;EV-ENTR03-04;EV-ENTR07-02;EV-ENTR08-03}} \\ \hline
Mockups relacionados & {\ttfamily\seqsplit{MU-15;MU-16;MU-40;MU-41}} \\ \hline
Disparador & Recepción solicita registrar un pago. \\ \hline
\end{tabularx}

\textbf{Precondiciones.}
\begin{itemize}
\item El usuario está autenticado y autorizado para registrar pagos.
\item La interfaz de registro de pagos está disponible.
\end{itemize}

\textbf{Flujo principal \texttt{CU-10-FP}.}
\begin{enumerate}
\item Recepción abre el registro de pago.
\item Ingresa monto, concepto, medio, referencia y fecha.
\item El sistema valida los datos de la operación.
\item El servicio de pagos registra el pago.
\item El sistema genera un comprobante interno asociado.
\item El sistema registra el evento requerido y devuelve el pago con su comprobante.
\end{enumerate}

\textbf{Flujo alternativo \texttt{CU-10-FA-01}.} \textit{Condición de disparo:} El medio de pago no requiere una referencia externa, como en un pago en efectivo. \textit{Apertura y retorno:} Se abre desde el paso 3 del flujo principal; después de registrar la referencia como no aplicable, retorna al paso 4 para registrar el pago.
\begin{enumerate}[label=\alph*.]
\item El sistema registra la referencia como no aplicable según la regla del medio.
\item La operación continúa con los demás datos obligatorios y genera el comprobante interno.
\end{enumerate}

\textbf{Excepción \texttt{CU-10-EX-01}.} \textit{Condición de disparo:} El monto es inválido o falta un dato obligatorio para el medio seleccionado.
\begin{enumerate}[label=\alph*.]
\item El sistema no registra el pago.
\item No se genera un comprobante de una operación inválida.
\end{enumerate}

\textbf{Poscondiciones.}
\begin{itemize}
\item El pago válido conserva monto, concepto, medio, referencia aplicable y fecha.
\item Existe un comprobante interno vinculado a la operación.
\end{itemize}

\subsubsection{CU-11 -- Registrar asistencia}
\begin{tabularx}{\textwidth}{|>{\bfseries}p{3.2cm}|X|}\hline
Actor principal & Recepción \\ \hline
Requisito relacionado & \texttt{RF-13} \\ \hline
Historia / criterio & \texttt{HU-11} / \texttt{CA-11} \\ \hline
Componente & Asistencia \\ \hline
Evidencia de origen & {\ttfamily\seqsplit{EV-ENTR02-03;EV-ENTR03-03;EV-ENTR04-02;EV-ENTR07-07;EV-ENTR08-02}} \\ \hline
Mockups relacionados & {\ttfamily\seqsplit{MU-20;MU-30;MU-46;MU-47}} \\ \hline
Disparador & Recepción solicita registrar la asistencia de un cliente. \\ \hline
\end{tabularx}

\textbf{Precondiciones.}
\begin{itemize}
\item El cliente existe.
\item El usuario tiene permiso para registrar asistencia.
\end{itemize}

\textbf{Flujo principal \texttt{CU-11-FP}.}
\begin{enumerate}
\item Recepción localiza al cliente.
\item El sistema consulta la vigencia de la membresía.
\item El sistema confirma que la vigencia permite el ingreso.
\item Recepción confirma el registro de asistencia.
\item El sistema crea un único evento de asistencia no biométrico.
\item El sistema devuelve la confirmación del evento.
\end{enumerate}

\textbf{Flujo alternativo \texttt{CU-11-FA-01}.} \textit{Condición de disparo:} La asistencia requiere una excepción operativa autorizada. \textit{Apertura y retorno:} Se abre desde el paso 3 del flujo principal; después de documentar la excepción autorizada, retorna al paso 4 para confirmar el registro de asistencia.
\begin{enumerate}[label=\alph*.]
\item El usuario autorizado documenta la excepción.
\item El sistema registra el evento de asistencia junto con la referencia de la excepción autorizada.
\end{enumerate}

\textbf{Excepción \texttt{CU-11-EX-01}.} \textit{Condición de disparo:} La membresía no es válida y no existe una excepción autorizada, o ya existe el mismo evento de asistencia.
\begin{enumerate}[label=\alph*.]
\item El sistema rechaza el registro.
\item No se crea un evento duplicado ni no autorizado.
\end{enumerate}

\textbf{Poscondiciones.}
\begin{itemize}
\item Existe como máximo un evento de asistencia válido para la operación.
\item El registro no incorpora datos biométricos.
\end{itemize}

\subsubsection{CU-12 -- Consultar asistencia}
\begin{tabularx}{\textwidth}{|>{\bfseries}p{3.2cm}|X|}\hline
Actor principal & Personal autorizado \\ \hline
Requisito relacionado & \texttt{RF-14} \\ \hline
Historia / criterio & \texttt{HU-12} / \texttt{CA-12} \\ \hline
Componente & Asistencia \\ \hline
Evidencia de origen & {\ttfamily\seqsplit{EV-ENTR03-01;EV-ENTR07-01;EV-ENTR08-05;EV-ENTR09-06}} \\ \hline
Mockups relacionados & {\ttfamily\seqsplit{MU-20;MU-46;MU-56;MU-65(F)}} \\ \hline
Disparador & Un usuario autorizado solicita consultar registros de asistencia. \\ \hline
\end{tabularx}

\textbf{Precondiciones.}
\begin{itemize}
\item El usuario mantiene una sesión válida.
\item Dispone de al menos un criterio permitido de consulta o acceso al resumen autorizado.
\end{itemize}

\textbf{Flujo principal \texttt{CU-12-FP}.}
\begin{enumerate}
\item El usuario abre la consulta de asistencia.
\item Selecciona filtros por cliente, fecha o turno.
\item El sistema valida los filtros.
\item El servicio de asistencia recupera los eventos autorizados.
\item El sistema calcula los conteos correspondientes.
\item El sistema muestra detalle y resumen de la consulta.
\end{enumerate}

\textbf{Flujo alternativo \texttt{CU-12-FA-01}.} \textit{Condición de disparo:} La combinación de filtros no devuelve eventos. \textit{Apertura y retorno:} Se abre desde el paso 4 del flujo principal; el caso de uso termina después de presentar el resultado vacío y los conteos en cero.
\begin{enumerate}[label=\alph*.]
\item El sistema presenta el resultado vacío con los filtros aplicados.
\item Los conteos se muestran en cero sin inventar registros.
\end{enumerate}

\textbf{Excepción \texttt{CU-12-EX-01}.} \textit{Condición de disparo:} El rango de fechas es inválido o el usuario intenta consultar información fuera de su permiso.
\begin{enumerate}[label=\alph*.]
\item El sistema no ejecuta la consulta solicitada.
\item No se expone información fuera del rol.
\end{enumerate}

\textbf{Poscondiciones.}
\begin{itemize}
\item La consulta no altera los eventos de asistencia.
\item El detalle y los conteos corresponden al mismo conjunto filtrado.
\end{itemize}

\subsubsection{CU-13 -- Administrar productos}
\begin{tabularx}{\textwidth}{|>{\bfseries}p{3.2cm}|X|}\hline
Actor principal & Administrador \\ \hline
Requisito relacionado & \texttt{RF-15} \\ \hline
Historia / criterio & \texttt{HU-13} / \texttt{CA-13} \\ \hline
Componente & Inventario \\ \hline
Evidencia de origen & {\ttfamily\seqsplit{EV-ENTR01-01;EV-ENTR02-05}} \\ \hline
Mockups relacionados & {\ttfamily\seqsplit{MU-19;MU-28;MU-44;MU-45}} \\ \hline
Disparador & El administrador solicita crear o mantener un producto. \\ \hline
\end{tabularx}

\textbf{Precondiciones.}
\begin{itemize}
\item El administrador está autenticado.
\item La gestión de productos está habilitada para su rol.
\end{itemize}

\textbf{Flujo principal \texttt{CU-13-FP}.}
\begin{enumerate}
\item El administrador abre el catálogo de productos.
\item Ingresa código, nombre, precio, unidad, stock mínimo y estado.
\item El sistema valida los datos.
\item El sistema comprueba la unicidad del código.
\item El sistema guarda el producto.
\item El catálogo refleja el registro actualizado.
\end{enumerate}

\textbf{Flujo alternativo \texttt{CU-13-FA-01}.} \textit{Condición de disparo:} El producto ya existe y el administrador solicita modificar sus datos permitidos o su estado. \textit{Apertura y retorno:} Se abre desde el paso 4 del flujo principal; después de actualizar el mismo registro, retorna al paso 6 para reflejar el producto actualizado en el catálogo.
\begin{enumerate}[label=\alph*.]
\item El sistema carga el producto existente.
\item El administrador modifica los datos permitidos y el sistema actualiza el mismo registro.
\end{enumerate}

\textbf{Excepción \texttt{CU-13-EX-01}.} \textit{Condición de disparo:} El código está duplicado o el precio, unidad o stock mínimo son inválidos.
\begin{enumerate}[label=\alph*.]
\item El sistema no crea ni actualiza el producto con datos inconsistentes.
\item El catálogo previo permanece sin cambios.
\end{enumerate}

\textbf{Poscondiciones.}
\begin{itemize}
\item El producto queda identificado por un código único.
\item El catálogo conserva precio, unidad, stock mínimo y estado consistentes.
\end{itemize}

\subsubsection{CU-14 -- Registrar movimientos de stock}
\begin{tabularx}{\textwidth}{|>{\bfseries}p{3.2cm}|X|}\hline
Actor principal & Personal autorizado \\ \hline
Requisito relacionado & \texttt{RF-16} \\ \hline
Historia / criterio & \texttt{HU-14} / \texttt{CA-14} \\ \hline
Componente & Inventario \\ \hline
Evidencia de origen & {\ttfamily\seqsplit{EV-ENTR01-04;EV-ENTR04-04;EV-ENTR07-03}} \\ \hline
Mockups relacionados & {\ttfamily\seqsplit{MU-69}} \\ \hline
Disparador & Un usuario autorizado solicita registrar una entrada o ajuste de inventario. \\ \hline
\end{tabularx}

\textbf{Precondiciones.}
\begin{itemize}
\item El usuario está autenticado y tiene permiso para registrar movimientos de stock.
\item La gestión de inventario está disponible para iniciar la operación.
\end{itemize}

\textbf{Flujo principal \texttt{CU-14-FP}.}
\begin{enumerate}
\item El usuario indica el producto sobre el que desea registrar el movimiento.
\item El sistema verifica que el producto exista en el catálogo.
\item El usuario selecciona una entrada e ingresa cantidad y motivo.
\item El sistema identifica al responsable desde la sesión y calcula el saldo resultante incrementando las existencias.
\item El sistema valida que la cantidad y el saldo resultante sean consistentes.
\item El sistema registra cantidad, motivo, responsable y saldo resultante.
\item El sistema confirma el movimiento y el nuevo balance.
\end{enumerate}

\textbf{Flujo alternativo \texttt{CU-14-FA-01}.} \textit{Condición de disparo:} En el paso 3 del flujo principal el usuario selecciona un ajuste en lugar de una entrada. \textit{Apertura y retorno:} Se abre desde el paso 3; después de calcular el efecto del ajuste, retorna al paso 5 para validar el saldo antes de registrar el movimiento.
\begin{enumerate}[label=\alph*.]
\item El usuario indica la cantidad y el motivo del ajuste, especificando su efecto sobre las existencias.
\item El sistema identifica al responsable desde la sesión y calcula el saldo resultante del ajuste.
\end{enumerate}

\textbf{Excepción \texttt{CU-14-EX-01}.} \textit{Condición de disparo:} El producto no existe, la cantidad es inválida o el movimiento produciría un saldo negativo.
\begin{enumerate}[label=\alph*.]
\item El sistema rechaza el movimiento.
\item El saldo anterior permanece sin cambios.
\end{enumerate}

\textbf{Poscondiciones.}
\begin{itemize}
\item El saldo resultante coincide con el movimiento registrado.
\item Cada movimiento conserva motivo y responsable.
\end{itemize}

\subsubsection{CU-15 -- Registrar venta}
\begin{tabularx}{\textwidth}{|>{\bfseries}p{3.2cm}|X|}\hline
Actor principal & Recepción \\ \hline
Requisito relacionado & \texttt{RF-17} \\ \hline
Historia / criterio & \texttt{HU-15} / \texttt{CA-15} \\ \hline
Componente & Ventas \\ \hline
Evidencia de origen & {\ttfamily\seqsplit{EV-ENTR02-05;EV-ENTR03-01;EV-ENTR04-04;EV-ENTR07-03;EV-ENTR08-03}} \\ \hline
Mockups relacionados & {\ttfamily\seqsplit{MU-18;MU-27;MU-42;MU-43}} \\ \hline
Disparador & Recepción solicita registrar una venta de productos. \\ \hline
\end{tabularx}

\textbf{Precondiciones.}
\begin{itemize}
\item El usuario está autorizado para ventas.
\item Los productos seleccionados existen.
\end{itemize}

\textbf{Flujo principal \texttt{CU-15-FP}.}
\begin{enumerate}
\item Recepción selecciona los productos y cantidades.
\item El sistema consulta el stock disponible.
\item El sistema valida que todas las cantidades tengan stock suficiente.
\item El sistema calcula la venta.
\item Recepción confirma la operación.
\item El sistema registra la venta y descuenta existencias en la misma operación demostrativa.
\item El sistema devuelve la venta, comprobante interno y saldos resultantes.
\end{enumerate}

\textbf{Flujo alternativo \texttt{CU-15-FA-01}.} \textit{Condición de disparo:} Uno de los ítems no dispone de la cantidad solicitada antes de confirmar la venta. \textit{Apertura y retorno:} Se abre desde el paso 3 del flujo principal; después de que Recepción ajusta o retira la cantidad, retorna al paso 2 para consultar nuevamente el stock y repetir la validación.
\begin{enumerate}[label=\alph*.]
\item El sistema identifica el ítem con stock insuficiente y solicita ajustar o retirar su cantidad.
\item Recepción corrige el detalle y el sistema vuelve a validar antes de permitir la confirmación.
\end{enumerate}

\textbf{Excepción \texttt{CU-15-EX-01}.} \textit{Condición de disparo:} Ocurre un fallo al persistir la venta o el descuento de stock.
\begin{enumerate}[label=\alph*.]
\item El sistema no deja una venta parcialmente registrada.
\item La venta y la actualización de existencias se revierten como una única operación lógica.
\end{enumerate}

\textbf{Poscondiciones.}
\begin{itemize}
\item Una venta confirmada y sus descuentos de stock son consistentes entre sí.
\item Una operación fallida no deja saldos parciales.
\end{itemize}

\subsubsection{CU-16 -- Cerrar turno y conciliar caja}
\begin{tabularx}{\textwidth}{|>{\bfseries}p{3.2cm}|X|}\hline
Actor principal & Recepción \\ \hline
Requisito relacionado & \texttt{RF-19} \\ \hline
Historia / criterio & \texttt{HU-16} / \texttt{CA-16} \\ \hline
Componente & Caja \\ \hline
Evidencia de origen & {\ttfamily\seqsplit{EV-ENTR02-01;EV-ENTR02-06;EV-ENTR03-05;EV-ENTR04-05;EV-ENTR07-03;EV-ENTR08-04}} \\ \hline
Mockups relacionados & {\ttfamily\seqsplit{MU-70}} \\ \hline
Disparador & Recepción solicita cerrar un turno abierto y conciliar caja. \\ \hline
\end{tabularx}

\textbf{Precondiciones.}
\begin{itemize}
\item El usuario está autenticado y autorizado para ejecutar el cierre.
\item La función de conciliación de caja está disponible para iniciar la operación.
\end{itemize}

\textbf{Flujo principal \texttt{CU-16-FP}.}
\begin{enumerate}
\item Recepción inicia el cierre del turno.
\item El sistema verifica que exista un turno abierto y recupera sus movimientos asociados.
\item El sistema consolida pagos y ventas del turno por medio.
\item Recepción ingresa el valor declarado de caja.
\item El sistema compara los valores calculados con el valor declarado y calcula la diferencia.
\item El sistema registra el cierre, los totales y la diferencia.
\item El sistema confirma el cierre del turno.
\end{enumerate}

\textbf{Flujo alternativo \texttt{CU-16-FA-01}.} \textit{Condición de disparo:} La diferencia calculada es distinta de cero. \textit{Apertura y retorno:} Se abre desde el paso 5 del flujo principal; después de dejar la diferencia identificada para seguimiento, retorna al paso 6 para registrar el cierre con dicha diferencia.
\begin{enumerate}[label=\alph*.]
\item El sistema marca la diferencia calculada como parte explícita de la conciliación.
\item La interfaz la presenta para su seguimiento operativo antes de registrar el cierre.
\end{enumerate}

\textbf{Excepción \texttt{CU-16-EX-01}.} \textit{Condición de disparo:} No existe un turno abierto o falta el valor requerido para la conciliación.
\begin{enumerate}[label=\alph*.]
\item El sistema impide cerrar el turno.
\item No se crea un registro de cierre incompleto.
\end{enumerate}

\textbf{Poscondiciones.}
\begin{itemize}
\item El turno queda cerrado con totales por medio y diferencia registrada.
\item Los movimientos previos permanecen trazables.
\end{itemize}

\subsubsection{CU-17 -- Gestionar novedades internas}
\begin{tabularx}{\textwidth}{|>{\bfseries}p{3.2cm}|X|}\hline
Actor principal & Personal autorizado \\ \hline
Requisito relacionado & \texttt{RF-20} \\ \hline
Historia / criterio & \texttt{HU-17} / \texttt{CA-17} \\ \hline
Componente & Novedades \\ \hline
Evidencia de origen & {\ttfamily\seqsplit{EV-ENTR02-01;EV-ENTR02-08;EV-ENTR03-07;EV-ENTR04-06;EV-ENTR05-01}} \\ \hline
Mockups relacionados & {\ttfamily\seqsplit{MU-21;MU-31;MU-50;MU-51;MU-61}} \\ \hline
Disparador & Un usuario autorizado solicita crear o actualizar una novedad interna. \\ \hline
\end{tabularx}

\textbf{Precondiciones.}
\begin{itemize}
\item El usuario mantiene una sesión válida.
\item La categoría de novedad está disponible para la operación.
\end{itemize}

\textbf{Flujo principal \texttt{CU-17-FP}.}
\begin{enumerate}
\item El usuario abre la gestión de novedades.
\item Ingresa categoría, detalle, responsable, vencimiento y estado.
\item El sistema valida los campos requeridos.
\item El sistema registra la novedad.
\item El sistema incorpora el estado inicial al historial.
\item El sistema confirma la novedad trazable.
\end{enumerate}

\textbf{Flujo alternativo \texttt{CU-17-FA-01}.} \textit{Condición de disparo:} Una novedad existente cambia de responsable, vencimiento o estado. \textit{Apertura y retorno:} Se abre desde el paso 1 del flujo principal cuando el usuario elige una novedad existente; el caso de uso termina después de registrar el cambio y conservar el estado anterior en el historial.
\begin{enumerate}[label=\alph*.]
\item El usuario selecciona la novedad y registra el cambio permitido.
\item El sistema conserva el estado anterior en el historial y añade el nuevo cambio.
\end{enumerate}

\textbf{Excepción \texttt{CU-17-EX-01}.} \textit{Condición de disparo:} Falta información obligatoria o el usuario no está autorizado para el cambio.
\begin{enumerate}[label=\alph*.]
\item El sistema rechaza la operación.
\item La novedad y su historial permanecen sin alteraciones.
\end{enumerate}

\textbf{Poscondiciones.}
\begin{itemize}
\item La novedad conserva sus datos operativos y el historial de estados.
\item Los cambios son atribuibles al usuario autorizado.
\end{itemize}

\subsubsection{CU-18 -- Crear y asignar rutina}
\begin{tabularx}{\textwidth}{|>{\bfseries}p{3.2cm}|X|}\hline
Actor principal & Instructor \\ \hline
Requisito relacionado & \texttt{RF-22} \\ \hline
Historia / criterio & \texttt{HU-18} / \texttt{CA-18} \\ \hline
Componente & Rutinas \\ \hline
Evidencia de origen & {\ttfamily\seqsplit{EV-ENTR01-05;EV-ENTR05-04;EV-ENTR09-01;EV-ENTR09-02;EV-ENTR09-03}} \\ \hline
Mockups relacionados & {\ttfamily\seqsplit{MU-17;MU-29;MU-54;MU-55;MU-64(F)}} \\ \hline
Disparador & El instructor solicita crear una rutina y asignarla a un cliente. \\ \hline
\end{tabularx}

\textbf{Precondiciones.}
\begin{itemize}
\item El instructor está autenticado y autorizado.
\item El cliente objetivo existe en el sistema.
\end{itemize}

\textbf{Flujo principal \texttt{CU-18-FP}.}
\begin{enumerate}
\item El instructor selecciona al cliente.
\item Define el objetivo operativo de la rutina.
\item Añade ejercicios con series, repeticiones, descanso y días.
\item El sistema valida la estructura requerida.
\item El instructor revisa y confirma la rutina.
\item El sistema crea la rutina y registra su asignación al cliente.
\item El sistema confirma la versión inicial sin incorporar datos de salud reales.
\end{enumerate}

\textbf{Flujo alternativo \texttt{CU-18-FA-01}.} \textit{Condición de disparo:} Antes de confirmar, el instructor modifica la selección u orden de ejercicios o sus parámetros. \textit{Apertura y retorno:} Se abre desde el paso 5 del flujo principal; después de actualizar la composición temporal, retorna al paso 4 para repetir la validación antes de confirmar.
\begin{enumerate}[label=\alph*.]
\item El sistema actualiza la composición temporal de la rutina.
\item La validación se repite y el instructor puede confirmar la versión ajustada.
\end{enumerate}

\textbf{Excepción \texttt{CU-18-EX-01}.} \textit{Condición de disparo:} Falta el objetivo o algún ejercicio carece de parámetros mínimos requeridos.
\begin{enumerate}[label=\alph*.]
\item El sistema no asigna la rutina incompleta.
\item Se solicita completar la información operativa sin pedir datos de salud excluidos.
\end{enumerate}

\textbf{Poscondiciones.}
\begin{itemize}
\item Existe una rutina asignada con objetivo, ejercicios y parámetros operativos.
\item La rutina no contiene datos de salud reales fuera del alcance.
\end{itemize}

\subsubsection{CU-19 -- Versionar y ajustar rutina}
\begin{tabularx}{\textwidth}{|>{\bfseries}p{3.2cm}|X|}\hline
Actor principal & Instructor \\ \hline
Requisito relacionado & \texttt{RF-23} \\ \hline
Historia / criterio & \texttt{HU-19} / \texttt{CA-19} \\ \hline
Componente & Rutinas \\ \hline
Evidencia de origen & {\ttfamily\seqsplit{EV-ENTR01-05;EV-ENTR09-01;EV-ENTR09-04;EV-ENTR10-01;EV-ENTR10-03}} \\ \hline
Mockups relacionados & {\ttfamily\seqsplit{MU-17;MU-54}} \\ \hline
Disparador & El instructor solicita modificar una rutina existente. \\ \hline
\end{tabularx}

\textbf{Precondiciones.}
\begin{itemize}
\item El instructor está autorizado.
\item Existe una rutina vigente con al menos una versión registrada.
\end{itemize}

\textbf{Flujo principal \texttt{CU-19-FP}.}
\begin{enumerate}
\item El instructor abre la rutina vigente.
\item Selecciona el ejercicio o parámetro que requiere ajuste.
\item Modifica el ejercicio, series, repeticiones, descanso o días según corresponda.
\item El sistema valida la nueva configuración.
\item El sistema crea una nueva versión de la rutina.
\item El sistema conserva las versiones anteriores en solo lectura.
\item El sistema confirma la nueva versión.
\end{enumerate}

\textbf{Flujo alternativo \texttt{CU-19-FA-01}.} \textit{Condición de disparo:} El ajuste modifica únicamente parámetros de una rutina sin sustituir sus ejercicios. \textit{Apertura y retorno:} Se abre desde el paso 3 del flujo principal; después de conservar la composición de ejercicios, retorna al paso 4 para validar la configuración y continuar con el versionado.
\begin{enumerate}[label=\alph*.]
\item El sistema conserva la composición de ejercicios.
\item El cambio de parámetros genera igualmente una nueva versión y mantiene el historial previo.
\end{enumerate}

\textbf{Excepción \texttt{CU-19-EX-01}.} \textit{Condición de disparo:} El usuario intenta editar directamente una versión histórica o la nueva configuración es inválida.
\begin{enumerate}[label=\alph*.]
\item El sistema bloquea la modificación directa de la versión histórica.
\item No se crea una nueva versión hasta que la configuración sea válida.
\end{enumerate}

\textbf{Poscondiciones.}
\begin{itemize}
\item La versión nueva queda identificada como vigente cuando la operación concluye.
\item Las versiones anteriores permanecen inmutables y consultables en solo lectura.
\end{itemize}

\section{Modelado del sistema}
El modelado de esta versión corresponde al conjunto UML definitivo del equipo y reemplaza las figuras de síntesis provisionales usadas durante las revisiones internas 2.1 y 2.2. La carpeta fuente se conserva en \path{03_Modelado/Diagramas_UML/}. La revisión de integración comprobó correspondencia nominal entre los 19 casos de uso Must, sus diagramas individuales y los 19 diagramas de secuencia; los RF Should permanecen especificados y trazados aunque no tengan CU individual obligatorio.

\begin{table}[H]\centering\small
\begin{tabularx}{\textwidth}{XrX}\toprule
Grupo & Cantidad & Alcance \\ \midrule
Contexto & 1 & Límites, actores y sistemas externos.\\
Casos de uso & 20 & 1 general + 19 individuales asociados a RF Must.\\
Actividades & 5 & Procesos operativos principales.\\
Secuencia & 19 & Interacciones de los 19 RF Must.\\
Estados & 4 & Membresía, pago, novedad y rutina.\\
Clases & 1 & Modelo de dominio y relaciones.\\
Componentes & 1 & Descomposición lógica de la solución.\\
Despliegue & 1 & Nodos y distribución objetivo.\\
i* & 2 & Dependencias estratégicas SD y razonamiento SR.\\ \midrule
\textbf{Total} & \textbf{54} & \textbf{Conjunto definitivo integrado.}\\
\bottomrule\end{tabularx}
\caption{Inventario del modelado definitivo incorporado a la ERS.}\label{tab:uml-inventario}
\end{table}

\clearpage
\begin{landscape}
\subsection{Contexto del sistema}
\begin{figure}[H]\centering
\includegraphics[width=.96\linewidth,height=.70\textheight,keepaspectratio]{01_Contexto/CTX_CONTEXTO_FabroGym.png}
\caption{Diagrama de contexto definitivo de FabroGym.}\label{fig:ctx-final}
\end{figure}
\end{landscape}

\clearpage
\begin{landscape}
\subsection{Modelado organizacional i*}
\begin{figure}[H]\centering
\includegraphics[width=.95\linewidth,height=.58\textheight,keepaspectratio]{09_iStar/01_SD/ISTAR_SD_FabroGym.png}
\caption{Modelo i* Strategic Dependency (SD) definitivo.}\label{fig:istar-sd-final}
\end{figure}
\end{landscape}
\clearpage
\begin{landscape}
\subsection*{Modelo i* Strategic Rationale}
\begin{figure}[H]\centering
\includegraphics[width=.95\linewidth,height=.62\textheight,keepaspectratio]{09_iStar/02_SR/ISTAR_SR_FabroGym.png}
\caption{Modelo i* Strategic Rationale (SR) definitivo.}\label{fig:istar-sr-final}
\end{figure}
\end{landscape}

\clearpage
\begin{landscape}
\subsection{Casos de uso}
\begin{figure}[H]\centering
\includegraphics[width=.98\linewidth,height=.68\textheight,keepaspectratio]{02_Casos_de_Uso/00_General/CU_GENERAL_FabroGym.png}
\caption{Diagrama general definitivo de casos de uso Must.}\label{fig:cu-general-final}
\end{figure}
\end{landscape}

\clearpage
\begin{landscape}
\subsection{Modelo de dominio}
\begin{figure}[H]\centering
\includegraphics[width=.98\linewidth,height=.68\textheight,keepaspectratio]{06_Clases/DC_CLASES_FabroGym.png}
\caption{Diagrama de clases definitivo de FabroGym.}\label{fig:clases-final}
\end{figure}
\end{landscape}

\clearpage
\begin{landscape}
\subsection{Componentes}
\begin{figure}[H]\centering
\includegraphics[width=.97\linewidth,height=.66\textheight,keepaspectratio]{07_Componentes/COMP_COMPONENTES_FabroGym.png}
\caption{Diagrama de componentes definitivo.}\label{fig:componentes-final}
\end{figure}
\end{landscape}

\clearpage
\subsection{Despliegue}
\begin{figure}[H]\centering
\includegraphics[width=.90\textwidth,height=.70\textheight,keepaspectratio]{08_Despliegue/DEP_DESPLIEGUE_FabroGym.png}
\caption{Diagrama de despliegue definitivo.}\label{fig:despliegue-final}
\end{figure}

\clearpage
\subsection{Máquina de estados principal}
\begin{figure}[H]\centering
\includegraphics[height=.66\textheight,width=.75\textwidth,keepaspectratio]{05_Estados/DE_01_Membresia.png}
\caption{Estados definitivos de la membresía.}\label{fig:estado-membresia-final}
\end{figure}

\section{Arquitectura, seguridad y datos}
\subsection{Arquitectura demostrativa}
El MVP estático separa interfaz y módulos lógicos en JavaScript y utiliza almacenamiento local para la demostración. La arquitectura objetivo de producción deberá sustituir el almacenamiento local por un repositorio persistente, incorporar autenticación segura de servidor y aplicar controles transaccionales y de auditoría antes del uso real.

\subsection{Privacidad desde el diseño}
La configuración por defecto minimiza datos y prohíbe salud, biometría, peso, medidas corporales, dirección e imagen corporal real. Las pruebas y demostraciones deben usar datos sintéticos. Las evidencias identificables se mantienen en la zona restringida cifrada y no forman parte del repositorio público ni del depósito FAIR \cite{lopdp2021,spdp2025}.

\subsection{Seguridad y sesiones}
El diseño terminal exige control por rol, denegación por defecto, expiración de sesión, protección de secretos y TLS en un despliegue conectado. Estos objetivos se verifican por RNF-09, RNF-10 y RNF-11; no se confunden con las credenciales académicas simples del prototipo demostrativo \cite{owasp2025,nist2025}.

\subsection{Accesibilidad y compatibilidad}
Las pantallas Must deberán cumplir criterios aplicables de WCAG 2.2 AA y funcionar en las dos versiones estables recientes de Chrome, Edge y Firefox al momento de la prueba \cite{wcag2024}.

\section{Trazabilidad terminal}
La matriz externa \texttt{04\_Trazabilidad/matriz\_trazabilidad.csv} contiene 162 filas: 97 trazas preexistentes migradas, 8 planes de verificación asociados al componente inteligente y 57 trazas explícitas de flujo para los 19 casos de uso Must. La columna \texttt{ID\_Flujo} relaciona cada \texttt{CU-xx-FP}, \texttt{CU-xx-FA-01} y \texttt{CU-xx-EX-01} con su requisito, historia, criterio de aceptación, evidencia, componente y mockups. La figura resume el encadenamiento utilizado.
\begin{figure}[H]\centering\includegraphics[width=.95\textwidth]{figuras_2B/fig_trazabilidad_2B.pdf}\caption{Cadena de trazabilidad de extremo a extremo.}\label{fig:traza}\end{figure}

\begin{landscape}\scriptsize
\begin{longtable}{p{2.2cm}p{4.0cm}p{4.0cm}p{2.2cm}p{2.2cm}p{2.2cm}p{3.5cm}}
\toprule RF & Evidencias & Stakeholder & CU & HU & CA & Mockups \\ \midrule\endfirsthead
\toprule RF & Evidencias & Stakeholder & CU & HU & CA & Mockups \\ \midrule\endhead
RF-01 & EV-ENTR02-07 & Personal autorizado & CU-01 & HU-01 & CA-01 & MU-01 \\
RF-02 & EV-ENTR02-07;EV-ENTR03-06;EV-ENTR08-04 & Administrador & CU-02 & HU-02 & CA-02 & MU-22;MU-32 \\
RF-03 & EV-ENTR02-04;EV-ENTR02-07;EV-ENTR03-06 & Administrador &  &  & CA-N/A & MU-25 \\
RF-04 & EV-ENTR01-02;EV-ENTR02-02;EV-ENTR03-02;EV-ENTR08-01 & Recepción & CU-03 & HU-03 & CA-03 & MU-12;MU-37 \\
RF-05 & EV-ENTR02-03;EV-ENTR02-08;EV-ENTR08-05 & Recepción & CU-04 & HU-04 & CA-04 & MU-11;MU-36 \\
RF-06 & EV-ENTR02-02;EV-ENTR03-02;EV-ENTR04-01;EV-ENTR05-02 & Recepción & CU-05 & HU-05 & CA-05 & MU-11;MU-36 \\
RF-07 & EV-ENTR02-01;EV-ENTR04-01 & Administrador & CU-06 & HU-06 & CA-06 & MU-68 \\
RF-08 & EV-ENTR01-02;EV-ENTR02-04;EV-ENTR03-01;EV-ENTR04-03 & Recepción & CU-07 & HU-07 & CA-07 & MU-14;MU-39 \\
RF-09 & EV-ENTR01-03;EV-ENTR02-03;EV-ENTR03-03;EV-ENTR08-02 & Personal autorizado & CU-08 & HU-08 & CA-08 & MU-13;MU-38;MU-63(F) \\
RF-10 & EV-ENTR01-03;EV-ENTR05-02;EV-ENTR07-02 & Recepción & CU-09 & HU-09 & CA-09 & MU-10;MU-13;MU-38 \\
RF-11 & EV-ENTR01-02;EV-ENTR02-04;EV-ENTR03-04;EV-ENTR07-02;EV-ENTR08-03 & Recepción & CU-10 & HU-10 & CA-10 & MU-15;MU-16;MU-40;MU-41 \\
RF-12 & EV-ENTR02-04;EV-ENTR03-04;EV-ENTR07-07 & Administrador &  &  & CA-N/A & MU-15;MU-40 \\
RF-13 & EV-ENTR02-03;EV-ENTR03-03;EV-ENTR04-02;EV-ENTR07-07;EV-ENTR08-02 & Recepción & CU-11 & HU-11 & CA-11 & MU-20;MU-30;MU-46;MU-47 \\
RF-14 & EV-ENTR03-01;EV-ENTR07-01;EV-ENTR08-05;EV-ENTR09-06 & Personal autorizado & CU-12 & HU-12 & CA-12 & MU-20;MU-46;MU-56;MU-65(F) \\
RF-15 & EV-ENTR01-01;EV-ENTR02-05 & Administrador & CU-13 & HU-13 & CA-13 & MU-19;MU-28;MU-44;MU-45 \\
RF-16 & EV-ENTR01-04;EV-ENTR04-04;EV-ENTR07-03 & Personal autorizado & CU-14 & HU-14 & CA-14 & MU-69 \\
RF-17 & EV-ENTR02-05;EV-ENTR03-01;EV-ENTR04-04;EV-ENTR07-03;EV-ENTR08-03 & Recepción & CU-15 & HU-15 & CA-15 & MU-18;MU-27;MU-42;MU-43 \\
RF-18 & EV-ENTR01-04;EV-ENTR02-05;EV-ENTR05-04;EV-ENTR06-04 & Administrador &  &  & CA-N/A & MU-10;MU-19;MU-44 \\
RF-19 & EV-ENTR02-01;EV-ENTR02-06;EV-ENTR03-05;EV-ENTR04-05;EV-ENTR07-03;EV-ENTR08-04 & Recepción & CU-16 & HU-16 & CA-16 & MU-70 \\
RF-20 & EV-ENTR02-01;EV-ENTR02-08;EV-ENTR03-07;EV-ENTR04-06;EV-ENTR05-01 & Personal autorizado & CU-17 & HU-17 & CA-17 & MU-21;MU-31;MU-50;MU-51;MU-61 \\
RF-21 & EV-ENTR07-06;EV-ENTR09-03;EV-ENTR10-02 & Administrador &  &  & CA-N/A & MU-23;MU-33;MU-48;MU-49;MU-58 \\
RF-22 & EV-ENTR01-05;EV-ENTR05-04;EV-ENTR09-01;EV-ENTR09-02;EV-ENTR09-03 & Instructor & CU-18 & HU-18 & CA-18 & MU-17;MU-29;MU-54;MU-55;MU-64(F) \\
RF-23 & EV-ENTR01-05;EV-ENTR09-01;EV-ENTR09-04;EV-ENTR10-01;EV-ENTR10-03 & Instructor & CU-19 & HU-19 & CA-19 & MU-17;MU-54 \\
RF-24 & EV-ENTR09-04;EV-ENTR09-05;EV-ENTR09-06;EV-ENTR10-03;EV-ENTR10-04 & Instructor &  &  & CA-N/A & MU-59;MU-60 \\
RF-25 & EV-ENTR01-05;EV-ENTR05-05;EV-ENTR06-04;EV-ENTR07-04;EV-ENTR08-05;EV-ENTR09-05 & Administrador &  &  & CA-N/A & MU-10;MU-24;MU-34 \\
\bottomrule\end{longtable}\end{landscape}

\section{Plan de verificación y aceptación}
\subsection{Tipos de verificación}
\begin{tabularx}{\textwidth}{p{3.7cm}p{3.6cm}X}\toprule
Tipo & Aplicación & Evidencia esperada \\ \midrule
Prueba funcional & RF Must y Should prototipados & pasos, entradas, resultado esperado/obtenido, captura o registro.\\
Inspección & privacidad, alcance, minimización & revisión de formularios, campos y repositorio público.\\
Prueba de rendimiento & RNF-02 y RNF-03 & dataset sintético, tiempos, percentil y tasa de error.\\
Walkthrough & RNF-05, RNF-06 y RNF-16 a RNF-19 & acta/transcripción anonimizada y codificación.\\
Auditoría técnica & seguridad, compatibilidad y accesibilidad & checklist, navegador, herramienta y versión.\\
Restauración/simulacro & fiabilidad y respaldo & bitácora de prueba y tiempo de recuperación.\\
\bottomrule\end{tabularx}

\subsection{Estado de pruebas}
La ERS especifica métodos y umbrales, pero no declara ejecutadas pruebas que no aparecen como salidas reproducibles o evidencia versionada. En particular, las pruebas de carga, cobertura automatizada de código, compatibilidad completa, WCAG y restauración deben conservar su evidencia cuando se ejecuten. Esta separación protege contra afirmar resultados inexistentes.

\section{Cierre empírico relevante para la ERS}
\subsection{Resultados de explicabilidad}
Los cuatro RNF de explicabilidad (RNF-16 a RNF-19) se incorporan porque cuentan con fragmentos de evidencia, operacionalización y decisión de member checking. El componente recomendado sigue PROPUESTO. No se calcula un porcentaje global de cobertura del marco de explicabilidad porque el instrumento no fija un universo cerrado y verificable de dimensiones; reportar un porcentaje habría creado un denominador no soportado.

El contraste de perfiles se documenta en \texttt{07\_Datos/resultados/tablas/tabla\_efecto\_perfiles.csv}. La fila principal usa seis sesiones WALK independientes, informa \texttt{n\_unidades=6} e \texttt{interpretable=NO} para inferencia poblacional, con delta de Cliff=0,555556 e IC95\% [-0,333333; 1,000000]. Por tanto, la mayor proporción observada en las sesiones técnicas se describe como patrón del caso y no como efecto generalizable.

\subsection{Desviaciones del prerregistro}
El registro OSF fue publicado el 29-08-2026. Los seis WALK ocurrieron antes de ese registro, por lo que se declaran como evidencia previa/formativa y se reporta la desviación. También se documentan la normalización de identificadores, la no aplicación de pruebas inferenciales no soportadas, la naturaleza descriptiva de la encuesta, el resultado de saturación y la ausencia de grabación del member checking. El archivo terminal es \texttt{06\_Experimento/osf\_deviations.pdf} \cite{nosek2018}.

\subsection{Reproducibilidad}
El análisis 2B dispone de una única entrada canónica en \texttt{07\_Datos/scripts/run\_all.py} y tablas/figuras regenerables desde los datos crudos versionados. \texttt{06\_Experimento/} conserva procedencia y artefactos históricos, pero no constituye una segunda cadena canónica. Toda cifra cuantitativa usada por la ERS debe mantenerse sincronizada con las salidas de \texttt{07\_Datos/} \cite{sandve2013,wilkinson2016}.

\section{Limitaciones y dependencias externas documentadas}
\begin{riskbox}
\textbf{Member checking sin grabación.} Existe evidencia documental y 12 decisiones reales, pero no audio/video. La ERS reconoce el riesgo frente a la redacción literal de la rúbrica y no inventa una grabación.
\end{riskbox}
\begin{riskbox}
\textbf{Saturación estricta.} La curva de códigos queda en 6,306\% para las últimas tres sesiones y no satisface literalmente el umbral de 5\%; la estabilización axial de 1,852\% es evidencia complementaria, no sustitución matemática.
\end{riskbox}
\begin{riskbox}
\textbf{Tamaño del efecto por perfil.} El delta de Cliff se calcula con seis unidades independientes (3 WALK técnicos y 3 no técnicos), no con 18 categorías derivadas. La salida marca \texttt{interpretable=NO} para inferencia poblacional: el IC95\% principal es amplio y el tamaño muestral por perfil es mínimo. Se reporta como descripción exploratoria y no como prueba de diferencia entre poblaciones.
\end{riskbox}
\begin{riskbox}
\textbf{DOI Zenodo.} El paquete final de replicación publicado está identificado mediante DOI Zenodo \href{https://doi.org/10.5281/zenodo.22237884}{10.5281/zenodo.22237884}. El identificador se incorporó de forma consistente en esta ERS, el README de publicación y el manuscrito fuente. La actualización de \texttt{CITATION.cff} permanece fuera de este bloque y no se declara resuelta aquí.
\end{riskbox}
\begin{riskbox}
\textbf{Modelado final del equipo.} Se integró el conjunto definitivo de 54 diagramas y se verificó su correspondencia nominal con los 19 RF Must, las entidades de dominio, componentes y despliegue. Los diagramas no modifican por sí mismos el alcance funcional: en caso de contradicción prevalecen los 52 requisitos terminales y la trazabilidad 2B.
\end{riskbox}

\appendix
\section{Resumen de la ficha técnica de las 16 sesiones}
La tabla conserva exactamente los identificadores terminales de sesión, la técnica, el perfil, la fecha y las duraciones declaradas. Los hashes de la ficha se preservan como metadatos; la coincidencia criptográfica con el multimedia real debe verificarse ejecutando SHA-256 sobre los archivos de la zona restringida.
\begin{landscape}
\scriptsize
\begin{longtable}{p{2.2cm}p{3.2cm}p{4.2cm}p{2.7cm}p{7.8cm}p{2.2cm}}
\toprule Código & Técnica & Perfil & Fecha & Archivo de video & Duración \\ \midrule\endfirsthead
\toprule Código & Técnica & Perfil & Fecha & Archivo de video & Duración \\ \midrule\endhead
ENTR-01 & Entrevista & Propietario/administracion & 2026-07-28 & 2026-07-28\_Video\_ENTR-01\_DuenoGym.mp4 & 00:11:09 \\
ENTR-02 & Entrevista & Recepcionista & 2026-07-27 & 2026-07-27\_Video\_ENTR-02\_RecepcionistaGym.mp4 & 00:22:13 \\
ENTR-03 & Entrevista & Recepcionista & 2026-07-27 & 2026-07-27\_Video\_ENTR-03\_RecepcionistaGym.mp4 & 00:22:17 \\
ENTR-04 & Entrevista & Recepcionista & 2026-07-27 & 2026-07-27\_Video\_ENTR-04\_RecepcionistaGym.mp4 & 00:24:05 \\
ENTR-05 & Entrevista & Recepcionista & 2026-07-28 & 2026-07-28\_Video\_ENTR-05\_RecepcionistaGym.mp4 & 00:20:10 \\
ENTR-06 & Entrevista & Propietario/administracion & 2026-07-29 & 2026-07-29\_Video\_ENTR-06\_DuenoGym.mp4 & 00:27:51 \\
ENTR-07 & Entrevista & Propietario/administracion & 2026-07-28 & 2026-07-28\_Video\_ENTR-07\_DuenoGym.mp4 & 00:37:56 \\
ENTR-08 & Entrevista & Recepcionista & 2026-07-28 & 2026-07-28\_Video\_ENTR-08\_RecepcionistaGym.mp4 & 00:17:41 \\
ENTR-09 & Entrevista & Instructor & 2026-07-28 & 2026-07-28\_Video\_ENTR-09\_InstructorGym.mp4 & 00:21:06 \\
ENTR-10 & Entrevista & Instructor & 2026-07-25 & 2026-07-25\_Video\_ENTR-10\_InstructorGym.mp4 & 00:18:07 \\
WALK-NTEC-01 & Walkthrough no tecnico & Instructor & 2026-08-16 & 2026-08-16\_Video\_WALK-NTEC-01\_Instructor.mp4 & 00:17:53 \\
WALK-NTEC-02 & Walkthrough no tecnico & Recepcionista & 2026-08-16 & 2026-08-16\_Video\_WALK-NTEC-02\_Recepcionista.mp4 & 00:17:57 \\
WALK-NTEC-03 & Walkthrough no tecnico & Propietario/administracion & 2026-08-22 & 2026-08-22\_Video\_WALK-NTEC-03\_DuenoGym.mp4 & 00:22:58 \\
WALK-TEC-01 & Walkthrough tecnico & Ingenieria de software & 2026-08-12 & 2026-08-12\_Video\_WALK-TEC-01\_Ing.Soft.mp4 & 00:24:46 \\
WALK-TEC-02 & Walkthrough tecnico & Ingenieria de software & 2026-08-12 & 2026-08-12\_Video\_WALK-TEC-02\_Ing.Soft.mp4 & 00:26:39 \\
WALK-TEC-03 & Walkthrough tecnico & Ingenieria de software & 2026-08-13 & 2026-08-13\_Video\_WALK-TEC-03\_Ing.Soft.mp4 & 00:45:20 \\
\bottomrule\end{longtable}\normalsize
\end{landscape}

\section{Resumen cuantitativo del cierre 2B}
\begin{tabularx}{\textwidth}{Xr}\toprule
Indicador & Resultado \\ \midrule
Requisitos funcionales & 25 \\
RF Must & 19 \\
RF Should & 6 \\
RNF de calidad & 15 \\
Restricciones de diseño & 4 \\
RNF de explicabilidad finales & 4 \\
Total de requisitos terminales en esta ERS & 52 \\
Filas de trazabilidad preexistentes / planes IA añadidos & 97 / 8 \\
Trazas de flujo de casos de uso añadidas & 57 \\
Sesiones acumuladas & 16 \\
Duración de videos & 06:18:08 \\
Respuestas de encuesta & 70 \\
Fragmentos WALK codificados & 76 \\
Códigos normalizados WALK & 37 \\
Categorías WALK & 18 \\
Fragmentos pertinentes de explicabilidad & 9 \\
Decisiones de member checking & 12 \\
RNF de explicabilidad derivados de evidencia & 4 \\
Promedio códigos nuevos últimas 3 sesiones / total & 6,306\% \\
Estabilización axial complementaria & 1,852\% \\
\bottomrule\end{tabularx}

\section{Criterios de calidad de esta especificación}
Antes del corte final, cada requisito debe conservar: identificador único; necesidad/fuente; redacción singular; verificabilidad; prioridad; estado; relación con evidencia; y, cuando sea Must, CU/HU/CA y representación de interfaz. Los RNF deben conservar métrica, umbral, condición y método. Los RNF-16 a RNF-19 deben conservar además su evidencia WALK y decisión de member checking; RNF-20 a RNF-23 deben conservar la trazabilidad explícita hacia métrica, umbral, método, responsable, frecuencia, diseño y caso de prueba.



\section{Anexo UML definitivo del equipo}
\begin{importantbox}
\textbf{Propósito del anexo.} Este anexo contiene el modelado detallado vigente que respalda la ERS 2B. A diferencia del anexo histórico de v2.1, estas figuras se consideran el conjunto definitivo entregado por el equipo. La semántica normativa continúa en los requisitos terminales; el UML los representa y facilita su defensa, pero no crea requisitos nuevos.
\end{importantbox}

\subsection{Casos de uso individuales Must}
\begin{figure}[p]\centering
\includegraphics[width=.96\textwidth,height=.72\textheight,keepaspectratio]{02_Casos_de_Uso/01_Individuales/DCU_01_CU-01_Autenticar_usuario.png}
\caption{CU-01: Autenticar usuario. Diagrama de caso de uso definitivo.}\label{fig:dcu-01}
\end{figure}
\begin{figure}[p]\centering
\includegraphics[width=.96\textwidth,height=.72\textheight,keepaspectratio]{02_Casos_de_Uso/01_Individuales/DCU_02_CU-02_Aplicar_permisos_por_rol.png}
\caption{CU-02: Aplicar permisos por rol. Diagrama de caso de uso definitivo.}\label{fig:dcu-02}
\end{figure}
\begin{figure}[p]\centering
\includegraphics[width=.96\textwidth,height=.72\textheight,keepaspectratio]{02_Casos_de_Uso/01_Individuales/DCU_03_CU-03_Registrar_cliente_minimo.png}
\caption{CU-03: Registrar cliente mínimo. Diagrama de caso de uso definitivo.}\label{fig:dcu-03}
\end{figure}
\begin{figure}[p]\centering
\includegraphics[width=.96\textwidth,height=.72\textheight,keepaspectratio]{02_Casos_de_Uso/01_Individuales/DCU_04_CU-04_Buscar_y_consultar_cliente.png}
\caption{CU-04: Buscar y consultar cliente. Diagrama de caso de uso definitivo.}\label{fig:dcu-04}
\end{figure}
\begin{figure}[p]\centering
\includegraphics[width=.96\textwidth,height=.72\textheight,keepaspectratio]{02_Casos_de_Uso/01_Individuales/DCU_05_CU-05_Actualizar_o_reactivar_cliente.png}
\caption{CU-05: Actualizar o reactivar cliente. Diagrama de caso de uso definitivo.}\label{fig:dcu-05}
\end{figure}
\begin{figure}[p]\centering
\includegraphics[width=.96\textwidth,height=.72\textheight,keepaspectratio]{02_Casos_de_Uso/01_Individuales/DCU_06_CU-06_Configurar_planes_y_promociones.png}
\caption{CU-06: Configurar planes y promociones. Diagrama de caso de uso definitivo.}\label{fig:dcu-06}
\end{figure}
\begin{figure}[p]\centering
\includegraphics[width=.96\textwidth,height=.72\textheight,keepaspectratio]{02_Casos_de_Uso/01_Individuales/DCU_07_CU-07_Activar_o_renovar_membresia.png}
\caption{CU-07: Activar o renovar membresía. Diagrama de caso de uso definitivo.}\label{fig:dcu-07}
\end{figure}
\begin{figure}[p]\centering
\includegraphics[width=.96\textwidth,height=.72\textheight,keepaspectratio]{02_Casos_de_Uso/01_Individuales/DCU_08_CU-08_Consultar_vigencia.png}
\caption{CU-08: Consultar vigencia. Diagrama de caso de uso definitivo.}\label{fig:dcu-08}
\end{figure}
\begin{figure}[p]\centering
\includegraphics[width=.96\textwidth,height=.72\textheight,keepaspectratio]{02_Casos_de_Uso/01_Individuales/DCU_09_CU-09_Alertar_vencimientos.png}
\caption{CU-09: Alertar vencimientos. Diagrama de caso de uso definitivo.}\label{fig:dcu-09}
\end{figure}
\begin{figure}[p]\centering
\includegraphics[width=.96\textwidth,height=.72\textheight,keepaspectratio]{02_Casos_de_Uso/01_Individuales/DCU_10_CU-10_Registrar_pago_y_comprobante.png}
\caption{CU-10: Registrar pago y comprobante. Diagrama de caso de uso definitivo.}\label{fig:dcu-10}
\end{figure}
\begin{figure}[p]\centering
\includegraphics[width=.96\textwidth,height=.72\textheight,keepaspectratio]{02_Casos_de_Uso/01_Individuales/DCU_11_CU-11_Registrar_asistencia.png}
\caption{CU-11: Registrar asistencia. Diagrama de caso de uso definitivo.}\label{fig:dcu-11}
\end{figure}
\begin{figure}[p]\centering
\includegraphics[width=.96\textwidth,height=.72\textheight,keepaspectratio]{02_Casos_de_Uso/01_Individuales/DCU_12_CU-12_Consultar_asistencia.png}
\caption{CU-12: Consultar asistencia. Diagrama de caso de uso definitivo.}\label{fig:dcu-12}
\end{figure}
\begin{figure}[p]\centering
\includegraphics[width=.96\textwidth,height=.72\textheight,keepaspectratio]{02_Casos_de_Uso/01_Individuales/DCU_13_CU-13_Administrar_productos.png}
\caption{CU-13: Administrar productos. Diagrama de caso de uso definitivo.}\label{fig:dcu-13}
\end{figure}
\begin{figure}[p]\centering
\includegraphics[width=.96\textwidth,height=.72\textheight,keepaspectratio]{02_Casos_de_Uso/01_Individuales/DCU_14_CU-14_Registrar_movimientos_de_stock.png}
\caption{CU-14: Registrar movimientos de stock. Diagrama de caso de uso definitivo.}\label{fig:dcu-14}
\end{figure}
\begin{figure}[p]\centering
\includegraphics[width=.96\textwidth,height=.72\textheight,keepaspectratio]{02_Casos_de_Uso/01_Individuales/DCU_15_CU-15_Registrar_venta.png}
\caption{CU-15: Registrar venta. Diagrama de caso de uso definitivo.}\label{fig:dcu-15}
\end{figure}
\begin{figure}[p]\centering
\includegraphics[width=.96\textwidth,height=.72\textheight,keepaspectratio]{02_Casos_de_Uso/01_Individuales/DCU_16_CU-16_Cerrar_turno_y_conciliar_caja.png}
\caption{CU-16: Cerrar turno y conciliar caja. Diagrama de caso de uso definitivo.}\label{fig:dcu-16}
\end{figure}
\begin{figure}[p]\centering
\includegraphics[width=.96\textwidth,height=.72\textheight,keepaspectratio]{02_Casos_de_Uso/01_Individuales/DCU_17_CU-17_Gestionar_novedades_internas.png}
\caption{CU-17: Gestionar novedades internas. Diagrama de caso de uso definitivo.}\label{fig:dcu-17}
\end{figure}
\begin{figure}[p]\centering
\includegraphics[width=.96\textwidth,height=.72\textheight,keepaspectratio]{02_Casos_de_Uso/01_Individuales/DCU_18_CU-18_Crear_y_asignar_rutina.png}
\caption{CU-18: Crear y asignar rutina. Diagrama de caso de uso definitivo.}\label{fig:dcu-18}
\end{figure}
\begin{figure}[p]\centering
\includegraphics[width=.96\textwidth,height=.72\textheight,keepaspectratio]{02_Casos_de_Uso/01_Individuales/DCU_19_CU-19_Versionar_y_ajustar_rutina.png}
\caption{CU-19: Versionar y ajustar rutina. Diagrama de caso de uso definitivo.}\label{fig:dcu-19}
\end{figure}

\clearpage
\subsection{Diagramas de secuencia de los 19 RF Must}
Los 19 diagramas de secuencia fueron reexportados desde la fuente editable de Visual Paradigm con los identificadores vigentes. Sus nombres de archivo, rótulos internos y leyendas emplean de manera coherente la numeración terminal RF-01 a RF-23 aplicable a los 19 requisitos Must.
\begin{landscape}\begin{figure}[p]\centering
\includegraphics[width=.96\linewidth,height=.76\textheight,keepaspectratio]{04_Secuencia/DS_01_RF-01.png}
\caption{Secuencia de RF-01: Autenticar usuario.}\label{fig:ds-01}
\end{figure}\end{landscape}
\begin{landscape}\begin{figure}[p]\centering
\includegraphics[width=.96\linewidth,height=.76\textheight,keepaspectratio]{04_Secuencia/DS_02_RF-02.png}
\caption{Secuencia de RF-02: Aplicar permisos por rol.}\label{fig:ds-02}
\end{figure}\end{landscape}
\begin{landscape}\begin{figure}[p]\centering
\includegraphics[width=.96\linewidth,height=.76\textheight,keepaspectratio]{04_Secuencia/DS_03_RF-04.png}
\caption{Secuencia de RF-04: Registrar cliente mínimo.}\label{fig:ds-03}
\end{figure}\end{landscape}
\begin{landscape}\begin{figure}[p]\centering
\includegraphics[width=.96\linewidth,height=.76\textheight,keepaspectratio]{04_Secuencia/DS_04_RF-05.png}
\caption{Secuencia de RF-05: Buscar y consultar cliente.}\label{fig:ds-04}
\end{figure}\end{landscape}
\begin{landscape}\begin{figure}[p]\centering
\includegraphics[width=.96\linewidth,height=.76\textheight,keepaspectratio]{04_Secuencia/DS_05_RF-06.png}
\caption{Secuencia de RF-06: Actualizar o reactivar cliente.}\label{fig:ds-05}
\end{figure}\end{landscape}
\begin{landscape}\begin{figure}[p]\centering
\includegraphics[width=.96\linewidth,height=.76\textheight,keepaspectratio]{04_Secuencia/DS_06_RF-07.png}
\caption{Secuencia de RF-07: Configurar planes y promociones.}\label{fig:ds-06}
\end{figure}\end{landscape}
\begin{landscape}\begin{figure}[p]\centering
\includegraphics[width=.96\linewidth,height=.76\textheight,keepaspectratio]{04_Secuencia/DS_07_RF-08.png}
\caption{Secuencia de RF-08: Activar o renovar membresía.}\label{fig:ds-07}
\end{figure}\end{landscape}
\begin{landscape}\begin{figure}[p]\centering
\includegraphics[width=.96\linewidth,height=.76\textheight,keepaspectratio]{04_Secuencia/DS_08_RF-09.png}
\caption{Secuencia de RF-09: Consultar vigencia.}\label{fig:ds-08}
\end{figure}\end{landscape}
\begin{landscape}\begin{figure}[p]\centering
\includegraphics[width=.96\linewidth,height=.76\textheight,keepaspectratio]{04_Secuencia/DS_09_RF-10.png}
\caption{Secuencia de RF-10: Alertar vencimientos.}\label{fig:ds-09}
\end{figure}\end{landscape}
\begin{landscape}\begin{figure}[p]\centering
\includegraphics[width=.96\linewidth,height=.76\textheight,keepaspectratio]{04_Secuencia/DS_10_RF-11.png}
\caption{Secuencia de RF-11: Registrar pago y comprobante.}\label{fig:ds-10}
\end{figure}\end{landscape}
\begin{landscape}\begin{figure}[p]\centering
\includegraphics[width=.96\linewidth,height=.76\textheight,keepaspectratio]{04_Secuencia/DS_11_RF-13.png}
\caption{Secuencia de RF-13: Registrar asistencia.}\label{fig:ds-11}
\end{figure}\end{landscape}
\begin{landscape}\begin{figure}[p]\centering
\includegraphics[width=.96\linewidth,height=.76\textheight,keepaspectratio]{04_Secuencia/DS_12_RF-14.png}
\caption{Secuencia de RF-14: Consultar asistencia.}\label{fig:ds-12}
\end{figure}\end{landscape}
\begin{landscape}\begin{figure}[p]\centering
\includegraphics[width=.96\linewidth,height=.76\textheight,keepaspectratio]{04_Secuencia/DS_13_RF-15.png}
\caption{Secuencia de RF-15: Administrar productos.}\label{fig:ds-13}
\end{figure}\end{landscape}
\begin{landscape}\begin{figure}[p]\centering
\includegraphics[width=.96\linewidth,height=.76\textheight,keepaspectratio]{04_Secuencia/DS_14_RF-16.png}
\caption{Secuencia de RF-16: Registrar movimientos de stock.}\label{fig:ds-14}
\end{figure}\end{landscape}
\begin{landscape}\begin{figure}[p]\centering
\includegraphics[width=.96\linewidth,height=.76\textheight,keepaspectratio]{04_Secuencia/DS_15_RF-17.png}
\caption{Secuencia de RF-17: Registrar venta.}\label{fig:ds-15}
\end{figure}\end{landscape}
\begin{landscape}\begin{figure}[p]\centering
\includegraphics[width=.96\linewidth,height=.76\textheight,keepaspectratio]{04_Secuencia/DS_16_RF-19.png}
\caption{Secuencia de RF-19: Cerrar turno y conciliar caja.}\label{fig:ds-16}
\end{figure}\end{landscape}
\begin{landscape}\begin{figure}[p]\centering
\includegraphics[width=.96\linewidth,height=.76\textheight,keepaspectratio]{04_Secuencia/DS_17_RF-20.png}
\caption{Secuencia de RF-20: Gestionar novedades internas.}\label{fig:ds-17}
\end{figure}\end{landscape}
\begin{landscape}\begin{figure}[p]\centering
\includegraphics[width=.96\linewidth,height=.76\textheight,keepaspectratio]{04_Secuencia/DS_18_RF-22.png}
\caption{Secuencia de RF-22: Crear y asignar rutina.}\label{fig:ds-18}
\end{figure}\end{landscape}
\begin{landscape}\begin{figure}[p]\centering
\includegraphics[width=.96\linewidth,height=.76\textheight,keepaspectratio]{04_Secuencia/DS_19_RF-23.png}
\caption{Secuencia de RF-23: Versionar y ajustar rutina.}\label{fig:ds-19}
\end{figure}\end{landscape}

\clearpage
\subsection{Diagramas de actividad}
\begin{figure}[H]\centering
\includegraphics[height=.74\textheight,width=.82\textwidth,keepaspectratio]{03_Actividades/DA_01_Alta_Renovacion.png}
\caption{Diagrama de actividad: Alta o renovación de membresía.}\label{fig:da-01}
\end{figure}
\begin{figure}[H]\centering
\includegraphics[height=.74\textheight,width=.82\textwidth,keepaspectratio]{03_Actividades/DA_02_Ingreso_Asistencia.png}
\caption{Diagrama de actividad: Ingreso y asistencia.}\label{fig:da-02}
\end{figure}
\begin{figure}[H]\centering
\includegraphics[height=.74\textheight,width=.82\textwidth,keepaspectratio]{03_Actividades/DA_03_Venta_Caja.png}
\caption{Diagrama de actividad: Venta y caja.}\label{fig:da-03}
\end{figure}
\begin{figure}[H]\centering
\includegraphics[height=.74\textheight,width=.82\textwidth,keepaspectratio]{03_Actividades/DA_04_Rutina_Seguimiento.png}
\caption{Diagrama de actividad: Rutina y seguimiento.}\label{fig:da-04}
\end{figure}
\begin{figure}[H]\centering
\includegraphics[height=.74\textheight,width=.82\textwidth,keepaspectratio]{03_Actividades/DA_05_Novedad_Turno.png}
\caption{Diagrama de actividad: Novedad y turno.}\label{fig:da-05}
\end{figure}

\clearpage
\subsection{Máquinas de estados}
\begin{figure}[H]\centering
\includegraphics[height=.72\textheight,width=.74\textwidth,keepaspectratio]{05_Estados/DE_01_Membresia.png}
\caption{Máquina de estados: Membresía.}\label{fig:de-01}
\end{figure}
\begin{figure}[H]\centering
\includegraphics[height=.72\textheight,width=.74\textwidth,keepaspectratio]{05_Estados/DE_02_Pago.png}
\caption{Máquina de estados: Pago.}\label{fig:de-02}
\end{figure}
\begin{figure}[H]\centering
\includegraphics[height=.72\textheight,width=.74\textwidth,keepaspectratio]{05_Estados/DE_03_Novedad.png}
\caption{Máquina de estados: Novedad.}\label{fig:de-03}
\end{figure}
\begin{figure}[H]\centering
\includegraphics[height=.72\textheight,width=.74\textwidth,keepaspectratio]{05_Estados/DE_04_Rutina.png}
\caption{Máquina de estados: Rutina.}\label{fig:de-04}
\end{figure}

\clearpage
\begin{thebibliography}{17}
\bibitem{iso29148} ISO/IEC/IEEE, \emph{ISO/IEC/IEEE 29148:2018 Systems and software engineering -- Life cycle processes -- Requirements engineering}, 2018.
\bibitem{iso25010} ISO/IEC, \emph{ISO/IEC 25010:2023 Systems and software engineering -- Systems and software Quality Requirements and Evaluation (SQuaRE) -- Product quality model}, 2023.
\bibitem{pohl2010} K. Pohl, \emph{Requirements Engineering: Fundamentals, Principles, and Techniques}. Springer, 2010.
\bibitem{swebok2024} H. Washizaki and IEEE Computer Society, \emph{Guide to the Software Engineering Body of Knowledge (SWEBOK), Version 4.0}, 2024.
\bibitem{birt2016} L. Birt, S. Scott, D. Cavers, C. Campbell, and F. Walter, ``Member Checking: A Tool to Enhance Trustworthiness or Merely a Nod to Validation?'' \emph{Qualitative Health Research}, vol. 26, no. 13, pp. 1802--1811, 2016.
\bibitem{hennink2017} M. M. Hennink, B. N. Kaiser, and V. C. Marconi, ``Code Saturation Versus Meaning Saturation: How Many Interviews Are Enough?'' \emph{Qualitative Health Research}, vol. 27, no. 4, pp. 591--608, 2017.
\bibitem{molleri2020} J. S. Molleri, K. Petersen, and E. Mendes, ``An empirically evaluated checklist for surveys in software engineering,'' \emph{Information and Software Technology}, vol. 119, p. 106240, 2020.
\bibitem{cohen1988} J. Cohen, \emph{Statistical Power Analysis for the Behavioral Sciences}, 2nd ed. Lawrence Erlbaum Associates, 1988.
\bibitem{cliff1993} N. Cliff, ``Dominance Statistics: Ordinal Analyses to Answer Ordinal Questions,'' \emph{Psychological Bulletin}, vol. 114, no. 3, pp. 494--509, 1993, doi: 10.1037/0033-2909.114.3.494.
\bibitem{lopdp2021} Asamblea Nacional del Ecuador, \emph{Ley Orgánica de Protección de Datos Personales}, Registro Oficial, Quinto Suplemento 459, 2021.
\bibitem{spdp2025} Superintendencia de Protección de Datos Personales, \emph{Guía de Protección de Datos Personales desde el Diseño y por Defecto}, 2025.
\bibitem{owasp2025} OWASP Foundation, \emph{Application Security Verification Standard 5.0.0}, 2025.
\bibitem{nist2025} D. Temoshok et al., \emph{NIST SP 800-63B-4 Digital Identity Guidelines: Authentication and Authenticator Management}. NIST, 2025.
\bibitem{wcag2024} W3C, \emph{Web Content Accessibility Guidelines (WCAG) 2.2}, 2024. [Online]. Available: https://www.w3.org/TR/WCAG22/
\bibitem{nosek2018} B. A. Nosek, C. R. Ebersole, A. C. DeHaven, and D. T. Mellor, ``The preregistration revolution,'' \emph{Proceedings of the National Academy of Sciences}, vol. 115, no. 11, pp. 2600--2606, 2018.
\bibitem{sandve2013} G. K. Sandve, A. Nekrutenko, J. Taylor, and E. Hovig, ``Ten simple rules for reproducible computational research,'' \emph{PLoS Computational Biology}, vol. 9, no. 10, p. e1003285, 2013.
\bibitem{wilkinson2016} M. D. Wilkinson et al., ``The FAIR Guiding Principles for scientific data management and stewardship,'' \emph{Scientific Data}, vol. 3, p. 160018, 2016.
\end{thebibliography}
\end{document}
